% =====================================================================
% Paper 42 (standalone): Tomita-Takesaki modular structure on truncated
% SU(2) spectral triples; the four-witness 2pi period at finite cutoff
% (Riemannian). Demoted 2026-09-28 (/qa trunk FULL, PI-approved): the
% finite closure is automatic for the posited integer-spectrum K = 2 J_z;
% the finite-cutoff identification reading is withdrawn.
% 2026-09-28 (PI decision): Paper 43 (Krein-carrier extension) archived;
% its surviving content is absorbed as Sec. sec:lorentzian_closure.
%
% Math.OA / J. Geom. Phys. style; companion to Paper 38 (SU(2)
% propinquity convergence on Camporesi-Higuchi truncated triples) and
% Paper 40 (universal 4/pi rate constant on compact Lie groups).
%
% Status: first draft, May 2026, post-Sprint-L0 / L1 / L1-tighten
% closure (2026-05-16). Built from internal sprint memos:
%   debug/lorentzian_l0_audit_memo.md
%   debug/lorentzian_partition_transfer_memo.md
%   debug/nieuviarts_scoping_memo.md
%   debug/l1_modular_hamiltonian_architecture_memo.md
%   debug/l1_infrastructure_audit_memo.md
%   debug/l1_witness_spec_memo.md
%   debug/l1_modular_hamiltonian_results_memo.md
%   debug/l1_tighten_tomita_results_memo.md
% =====================================================================
\documentclass[11pt,a4paper]{article}

\usepackage[T1]{fontenc}
\usepackage[utf8]{inputenc}
\usepackage{amsmath, amssymb, amsthm, mathrsfs, mathtools}
\usepackage{geometry}
\geometry{margin=1in}
\usepackage{hyperref}
\hypersetup{colorlinks=true, linkcolor=blue!50!black, citecolor=blue!50!black, urlcolor=blue!50!black}
% \usepackage{microtype}  % disabled: MiKTeX font-expansion environmental issue
\usepackage{graphicx}
\usepackage{booktabs}
\usepackage[numbers,sort&compress]{natbib}
\usepackage{xcolor}

\theoremstyle{plain}
\newtheorem{theorem}{Theorem}[section]
\newtheorem{lemma}[theorem]{Lemma}
\newtheorem{proposition}[theorem]{Proposition}
\newtheorem{corollary}[theorem]{Corollary}
\theoremstyle{definition}
\newtheorem{definition}[theorem]{Definition}
\newtheorem{remark}[theorem]{Remark}
\newtheorem{example}[theorem]{Example}

\newcommand{\nmax}{n_{\max}}
\newcommand{\sthree}{S^{3}}
\newcommand{\Hilb}{\mathcal{H}}
\newcommand{\Op}{\mathcal{O}}
\newcommand{\Tcal}{\mathcal{T}}
\newcommand{\Acal}{\mathcal{A}}
\newcommand{\opnorm}[1]{\lVert #1 \rVert_{\mathrm{op}}}
\newcommand{\Frob}[1]{\lVert #1 \rVert_{F}}
\newcommand{\norm}[1]{\lVert #1 \rVert}
\newcommand{\R}{\mathbb{R}}
\newcommand{\C}{\mathbb{C}}
\newcommand{\T}{\mathbb{T}}
\newcommand{\N}{\mathbb{N}}
\newcommand{\Z}{\mathbb{Z}}
\newcommand{\Q}{\mathbb{Q}}
\newcommand{\SU}{\mathrm{SU}}
\newcommand{\SO}{\mathrm{SO}}
\newcommand{\Vol}{\mathrm{Vol}}
\newcommand{\Tr}{\mathrm{Tr}}
\newcommand{\Mat}{M}
\newcommand{\DCH}{D_{\mathrm{CH}}}
\newcommand{\KMS}{\mathrm{KMS}}
\newcommand{\TT}{\mathrm{TT}}
\newcommand{\BW}{\mathrm{BW}}
\newcommand{\HH}{\mathrm{HH}}
\newcommand{\Sew}{\mathrm{Sew}}
\newcommand{\Unruh}{\mathrm{Unruh}}
\newcommand{\GNS}{\mathrm{GNS}}
\newcommand{\Mod}{\mathrm{mod}}
\newcommand{\Bcal}{\mathcal{B}}
\newcommand{\Krein}{\mathcal{K}}
\newcommand{\Nt}{N_t}
\newcommand{\DGV}{D_{\mathrm{GV}}}
\newcommand{\DL}{D_{L}}
\newcommand{\KL}{K_{L}}
\newcommand{\JL}{J_{L}}
\newcommand{\Cl}{\mathrm{Cl}}

\title{Tomita--Takesaki modular structure on truncated\\
$\SU(2)$ spectral triples:\ the four-witness $2\pi$ period\\
at finite cutoff}
\author{J.~Loutey\thanks{Independent researcher. \texttt{jloutey@gmail.com}.
This work was developed in an AI-augmented agentic workflow with
Anthropic's Claude. Implementation, internal memos, and bibliographic
collation were performed collaboratively; all theorems, design choices,
and editorial decisions are the author's responsibility.}}
\date{May 2026 (preprint)}

\begin{document}
\maketitle

\begin{abstract}
The four-witness Wick-rotation theorem
(Hartle--Hawking~\cite{hartle_hawking1976}, Sewell~\cite{sewell1982},
Bisognano--Wichmann~\cite{bisognano_wichmann1976},
Unruh~\cite{unruh1976}) asserts that a state of a quantum field
restricted to a region bounded by a bifurcate Killing horizon with
surface gravity $\kappa_g$ is KMS at inverse temperature
$\beta = 2\pi/\kappa_g$ with respect to Killing-time evolution. The
shared $2\pi$ across all four faces is $\Vol(S^1)$, the regularity
period of the Euclidean rotation of the observer's causal-domain
coordinates. That continuum statement is prior art and is not
re-proved here. We study its finite-cutoff \emph{period} on the
Connes--van~Suijlekom~\cite{connes_vs2021} truncated
Camporesi--Higuchi~\cite{camporesi_higuchi1996} spectral triple
$\Tcal_{\nmax}$, with a hemispheric wedge and a posited generator
$K_\alpha^W = \mathrm{diag}(2m_j) = 2J_z$, the doubled rotation
generator about the polar axis. Because $2J_z$ has integer spectrum,
$e^{2\pi i K_\alpha^W} = I$, and the conjugation flow closes,
$\sigma_{2\pi}(O) = O$, for \emph{every} operator $O$ --- an elementary
fact [SYMBOLIC] that holds for any integer-spectrum generator
whatsoever and therefore does not single out a Bisognano--Wichmann
modular flow. The Tomita--Takesaki reading is not an independent
check:\ the wedge state is \emph{defined} as
$\rho_W = e^{-K_\alpha^W}/Z$, so Tomita's theorem for a Gibbs state
returns the same generator, $\sigma_t^{\TT} = \sigma_{-t}^{\alpha}$.
Nor is the agreement of the four witnesses (Bisognano--Wichmann,
Hartle--Hawking, Sewell, Unruh) a finite-cutoff result:\ the six
witness instantiations use the same $K_\alpha^W$ with the surface
gravity $\kappa_g$ relabelled, and $\beta H_{\mathrm{local}} =
K_\alpha^W$ is $\kappa_g$-independent by construction, so their
bit-identical residuals hold by construction. We record these closures
(bit-exact in floating point at $\nmax \in \{2, 3, 4, 5\}$, residual
$O(\sqrt{\dim \Hilb_{\nmax}} \cdot \varepsilon_{\mathrm{machine}})$) as
what they are. We also record the operator-level distinction between
$2J_z$ ($e^{2\pi i K} = +I$) and the wedge-restricted
Camporesi--Higuchi Dirac (half-integer spectrum, $e^{2\pi i D_W} = -I$,
the spinor double-cover sign;\ the conjugation flow closes for both).
An earlier version of this paper read these closures as lifting the
four-witness theorem to a ``literal identification'' at the % [retracted 2026-09-28: finite-four-witness-literal-identification]
operator-system level;\ that reading is withdrawn (2026-09-28). The
genuine finite-cutoff question --- \emph{derive} $K$ as the modular
Hamiltonian of a wedge-restricted Camporesi--Higuchi state, rather than
posit it --- is [OPEN]. The Lorentzian extension to signature $(3, 1)$ via
Bizi--Brouder--Besnard~2018~\cite{bizi_brouder_besnard2018}
Krein-space spectral triples is carried out \emph{at finite cutoff} in
the present paper (\S\ref{sec:lorentzian_closure}, which absorbs the
former Paper~43):\ a Krein carrier and Krein-self-adjoint Lorentzian
Dirac, an axiom audit at Bizi--Brouder--Besnard signature $(4, 6)$ in
which the universal chirality anticommutation fails on the truthful
Camporesi--Higuchi Dirac, the same automatic (and signature-blind) period
closure, and a Hilbert--Schmidt orthogonality between the posited local
Hamiltonian and the wedge Dirac that is a chirality-parity identity [SYMBOLIC],
the only transcendental of its residual, $1/\pi^2$, being the posited
$\beta = 2\pi/\kappa_g$ squared.  Its \emph{continuum}
(Lorentzian-propinquity-rate) closure was described here as the named
open frontier when this paper was drafted;\ that is
\textbf{superseded} (2026-06, Paper~45~\cite{paper45}):\ the natural
device fails structurally, and the chase is closed on the convention
branch (an [ARGUED] reading:\ its Euclidean-only leg is
generator-independent, its Wick-rotation/KMS-circle leg is conditional on
O0 for the Bisognano--Wichmann generator;\ Paper~45 Q1).  The mechanism is visible
in this paper's own result~(i):\ the posited wedge generator has
\emph{integer} spectrum, so the
finite-cutoff flow is a compact $\beta = 2\pi$ circle rather
than a non-compact boost [SYMBOLIC], and on the argued reading [ARGUED]
every Lipschitz/boost seminorm built on it is signature-blind.  The finite-cutoff closures reported below stand
as computed;\ what does not follow from them is a metric statement
about \emph{signature}, or an identification of the posited generator
with a Bisognano--Wichmann boost.
The construction is a finite-dimensional illustration of the
Connes--Rovelli thermal-time setup~\cite{connes_rovelli1994} on a
hemispheric wedge of the round $\sthree$, not a derivation of it.
\end{abstract}

\bigskip\noindent
\textbf{MSC2020 classifications:}\ 58B34 (noncommutative geometry),
46L60 (applications of selfadjoint operator algebras to physics),
46L87 (noncommutative Riemannian and metric geometry),
22E45 (representations of compact Lie groups),
81T05 (axiomatic quantum field theory).

\medskip\noindent
\textbf{Keywords:}\ Tomita--Takesaki modular flow, Bisognano--Wichmann
theorem, Connes--van~Suijlekom truncated operator system,
Camporesi--Higuchi Dirac, four-witness Wick rotation,
Hartle--Hawking, Sewell, Unruh effect, KMS state, $\SU(2)$ spectral
triple.

% =====================================================================
\section{Introduction}\label{sec:intro}
% =====================================================================

A central thread of axiomatic quantum field theory is the
\emph{four-witness Wick-rotation theorem}: under suitable analytic
hypotheses, a state of a quantum field restricted to a region
bounded by a bifurcate Killing horizon with surface gravity
$\kappa_g$ is a KMS state at inverse temperature
$\beta = 2\pi / \kappa_g$ with respect to evolution along the
relevant Killing vector field. The four physical instantiations are
Bisognano--Wichmann~\cite{bisognano_wichmann1976} for the Rindler
wedge of Minkowski spacetime, Hartle--Hawking~\cite{hartle_hawking1976}
for the Hartle--Hawking thermal vacuum on the Schwarzschild exterior,
Sewell~\cite{sewell1982} for the abstract KMS-modular reading of the
Hartle--Hawking vacuum on globally hyperbolic manifolds with bifurcate
Killing horizons, and Unruh~\cite{unruh1976} for the proper-acceleration
parameterisation of the Bisognano--Wichmann setup. The shared $2\pi$
across all four faces is $\Vol(S^1)$, the regularity period of the
Euclidean rotation of the observer's causal-domain coordinates after
Wick rotation $\tau = i t$.

The theorem closes naturally in axiomatic Wightman QFT for free fields
on globally hyperbolic Lorentzian spacetimes. Its operator-system
analog on a finite-dimensional truncation of a spectral triple --- i.e.,
the question whether the modular automorphism associated to a
wedge-restricted Gibbs state on a finite-dimensional operator system
realises the same $\sigma_{i \cdot 2\pi} = $ identity period closure
intrinsically, without recourse to the Wightman axiomatic Wick-rotation
chain --- is structurally distinct, and has not been settled in the
spectral-truncation literature
of~\cite{connes_vs2021, leimbach_vs2024}.

The present paper does \emph{not} settle that question.  It sets up
the finite-cutoff ingredients (wedge, Gibbs state, GNS/Tomita
construction) on the truncated Camporesi--Higuchi spectral triple
$\Tcal_{\nmax}$ over the round $\sthree = \SU(2)$, records what holds
when the modular generator is \emph{posited} to be the doubled
rotation generator $2J_z$, and isolates what remains open:\ deriving
the generator from a wedge-restricted Camporesi--Higuchi state.
We work in the setup of
Connes--van~Suijlekom~\cite{connes_vs2021}, with the truncated
operator system $\Op_{\nmax} = P_{\nmax} \, C^\infty(\sthree) \,
P_{\nmax}$ obtained by orthogonal projection onto the
finite-dimensional subspace $\Hilb_{\nmax}$ spanned by the lowest
Camporesi--Higuchi spinor harmonics, and with the Connes--vS
``Toeplitz S$^1$'' propagation invariant
$\mathrm{prop}(\Op_{\nmax}) = 2$ verified verbatim in supporting work
of the present author (Paper~32~\S~III). The convergence of
$\Tcal_{\nmax}$ to the round-$\sthree$ Camporesi--Higuchi spectral
triple $\Tcal_{\sthree}$ in van~Suijlekom's state-space
Gromov--Hausdorff distance --- with explicit rate
$\Lambda(\Tcal_{\nmax}, \Tcal_{\sthree}) \le C_3 \cdot \gamma_{\nmax}$
where $\gamma_{\nmax} \sim (4/\pi) \log \nmax / \nmax$ and $C_3 = 1$
--- has been established in the present author's Paper~38~\cite{paper38}.

\begin{theorem}[Main result, {[SYMBOLIC]}; automatic for any
integer-spectrum generator; cf.~Section~\ref{sec:unified_closure}]
\label{thm:main_intro}
Let $\Tcal_{\nmax}$ denote the truncated Camporesi--Higuchi spectral
triple of cutoff $\nmax$ over the round $\sthree = \SU(2)$, let
$P_W$ denote the hemispheric wedge projector defined by the equatorial
reflection $R_{\mathrm{polar}}: m_j \mapsto -m_j$ on the full-Dirac
spinor basis, and let
$\omega_W^{\BW}(O) = Z^{-1} \Tr_{P_W \Hilb_{\nmax}}(e^{-K_\alpha^W}\,O)$
denote the wedge-restricted Gibbs state at the BW local Hamiltonian
$H_{\mathrm{local}} = K_\alpha^W / \beta$. Then for every $\nmax \ge
1$, both the geometric BW-$\alpha$ modular flow
$\sigma_t^{\alpha}(O) = e^{itK_\alpha^W} O e^{-itK_\alpha^W}$ and the
Tomita--Takesaki BW-$\gamma$ modular flow $\sigma_t^{\TT}(O) =
\Delta^{it} O \Delta^{-it}$, with $\Delta = S^* S$ extracted from the
polar decomposition $S = J_{\TT} \Delta^{1/2}$ on the GNS Hilbert--Schmidt
space, satisfy the period closure
\begin{equation}\label{eq:period_closure_main}
   \sigma_{2\pi}^{\alpha}(O) \;=\; O,\qquad
   \sigma_{2\pi}^{\TT}(O) \;=\; O
\end{equation}
bit-exactly for every $O \in P_W \Op_{\nmax} P_W$, with maximum
periodicity residual scaling as
$O(\sqrt{\dim \Hilb_{\nmax}} \cdot \varepsilon_{\mathrm{machine}})$.
Furthermore, the two flows are operator-action conjugates,
\begin{equation}\label{eq:tt_alpha_conjugate_main}
   \sigma_t^{\TT}(O) \;=\; \sigma_{-t}^{\alpha}(O)
\end{equation}
bit-exact at every real $t$, and the density matrix
$\rho_W = e^{-K_\alpha^W}/Z$ is $\beta$-independent, so the six
witness instantiations $(\BW, \HH_{M=1}, \HH_{M=2},
\Sew_{M=1}, \Unruh_{a=1}, \Unruh_{a=2})$ produce identical
$\Delta$ and $K_{\TT}$.
\end{theorem}

Every clause of Theorem~\ref{thm:main_intro} is elementary, and none
of it is evidence for a Bisognano--Wichmann identification at finite
cutoff.  (i)~The period closure holds for \emph{any} Hermitian
generator with integer spectrum and \emph{any} operator $O$:\
$e^{2\pi i n} = 1$ for integer $n$.  Replacing $2J_z$ by an unrelated
integer-spectrum matrix leaves every closure test green.
(ii)~The conjugacy $\sigma_t^{\TT} = \sigma_{-t}^{\alpha}$ is Tomita's
theorem for a Gibbs state $\rho = e^{-K}/Z$, which returns $K$;\ since
the state was built from $K_\alpha^W$, the Tomita construction is
circular as a check on $K_\alpha^W$.  (iii)~The witness agreement is
by construction:\ the six instantiations share $K_\alpha^W$ and differ
only in the label $\kappa_g$, which cancels in $\beta H_{\mathrm{local}}$.
The generator $K_\alpha^W = \mathrm{two}\_m_j$ (odd integers in
$\{-(2l+1), \ldots, +(2l+1)\}$) is a \emph{rotation} generator on the
Riemannian $\sthree$, chosen, not derived, to play the role of the
boost.  The continuum four-witness statement
($\beta = 2\pi/\kappa_g$;\
\cite{bisognano_wichmann1976, hartle_hawking1976, sewell1982, unruh1976})
is prior art and is unaffected;\ its finite-cutoff counterpart ---
obtaining $K$ as the modular Hamiltonian of a wedge-restricted
Camporesi--Higuchi state and \emph{then} checking the period --- is
[OPEN] (Section~\ref{sec:conclusion}, O0).

\paragraph{What this paper does and does not close.}
This paper does \emph{not} close the principal Sprint
L1\footnote{Sprint/track designations are anchors into the project's internal chronicle (\texttt{CHANGELOG.md} in the repository).} falsifier named in
Paper~32~\S~VIII Remark~\texttt{rem:bisognano\_wichmann\_reading} of
the present author's framework documentation:\
``compute the modular Hamiltonian $K$ on $\Tcal_{\nmax}$ for a
wedge-restricted Camporesi--Higuchi state and verify $\sigma_{i \cdot 2\pi}
= $ identity at finite $\nmax$.'' That falsifier asks for $K$ to be
\emph{computed from} a Camporesi--Higuchi state;\ here $K$ is posited
($2J_z$) and the state is built from it, so the period closure is
automatic and tests nothing about the state.  The supporting module
reports the labels STRONG\_IDENTIFICATION (BW-$\alpha$) and
UNIFIED\_STRONG (BW-$\gamma$);\ they record that the automatic closure
was reproduced numerically, and are \emph{not} an identification
verdict (withdrawn 2026-09-28).  The falsifier remains [OPEN].
The Lorentzian extension to signature $(3, 1)$ via the
Bizi--Brouder--Besnard~2018~\cite{bizi_brouder_besnard2018} Krein-space
spectral triple is carried out \emph{at finite cutoff} in
\S\ref{sec:lorentzian_closure} (the same automatic period closure
$\sigma_{2\pi} = \mathrm{id}$ on the Krein wedge, not a four-witness
identification); the
continuum / Lorentzian-propinquity-rate extension is \emph{not} closed
by this paper (the named open target, and Paper~45's degeneracy theorem
closes off the natural Lorentzian-propinquity route). Sprint L0 of
the present author's framework documentation
(\texttt{debug/lorentzian\_l0\_audit\_memo.md},
Paper~31~\S~8, Paper~34~\S~V.E) audits the 28 projections of the
framework's projection taxonomy and predicts that the M3
sub-mechanism of the master Mellin engine becomes trivially empty at
signature $(3, 1)$ on chirality-symmetric spectrum, leaving the M1
$\Vol(S^1) = 2\pi$ mechanism (the load-bearing engine for the
four-witness theorem) intact under signature change;\ that prediction
turns out to be convention-dependent (\S\ref{sec:m3_convention}). The automatic
closure of the present paper does not confirm an operator-system
identification at signature $(3, 0)$;\ it shows only that the
posited generator has the integer spectrum a $2\pi$ period requires.

\paragraph{Related and concurrent work.}
The most directly comparable construction in the published literature
is the Connes--Rovelli~\cite{connes_rovelli1994} thermal time
hypothesis, which identifies the Tomita modular flow on a covariant
quantum field theory as the physical time evolution. The
Connes--Rovelli identification is general-covariant and abstract;
the present paper specialises it to a finite-dimensional operator
system in the Connes--van~Suijlekom~\cite{connes_vs2021} truncation
framework on the round $\sthree$. The continuum side has the long
tradition starting from Bisognano--Wichmann~\cite{bisognano_wichmann1976}
and standardised through Casini--Huerta~\cite{casini_huerta2009} and
Casini--Huerta--Myers~\cite{casini_huerta_myers2011}, with the
continuum modular Hamiltonian for a hemisphere of $S^d$ in CFT given
explicitly as $K_{\mathrm{cont}} = 2\pi \int_{H_+} \xi \cdot T^{00}\,d\Omega$
where $\xi$ is the wedge-preserving Killing field. The present
finite-cutoff construction mimics only the \emph{form} of this
identification:\ the integer spectrum of the posited generator stands
in for the canonical $2\pi$-period of the continuum modular flow, but
the generator is not obtained from a state. A finite-lattice
realisation in critical quantum spin chains was given by
Zhang~\emph{et~al.}~\cite{zhu_casini2020}, where the BW form is
\emph{extracted} from a ground state and so is a genuine test; the
present construction does not attempt that extraction, and differs
also in that the underlying object is a continuum
spectral triple on a non-flat compact $\sthree$ background, and the
finite-dimensional structure comes from the spectral truncation of
$\DCH$ rather than from a discretisation of physical space. The
existing modular literature on Riemannian backgrounds includes
Strohmaier~\cite{strohmaier2006} (Reeh--Schlieder on stationary
spacetimes) and Verch~\cite{verch2001} (generally covariant
spin-statistics and modular structure on curved
backgrounds); these address the algebraic-QFT side, whereas the
present paper anchors the modular structure on the
Connes--van~Suijlekom truncated operator system. Direct attempts to
lift the Wick-rotation chain at the spectral-triple level via
twisted-spectral-triple morphisms have been pursued recently by
Nieuviarts~\cite{nieuviarts2025a, nieuviarts2024,
nieuviarts2025b_proceedings}, but the twist construction
(Definition~2.2 of~\cite{nieuviarts2025a}) is explicitly restricted to
even-dimensional manifolds and does not apply to the $\sthree$ case
treated here.

\paragraph{Outline.}
Section~\ref{sec:setup} recalls the truncated Camporesi--Higuchi
spectral triple and the Paper~38 state-space Gromov--Hausdorff convergence result.
Section~\ref{sec:lorentzian_readiness} summarises the Sprint L0
Lorentzian-readiness audit, names the Riemannian scope of the present
paper, and footnotes the Nieuviarts NO-GO on direct twist-morphism
emergence for odd-dim $\sthree$. Section~\ref{sec:bw_at_finite_cutoff}
sets up the hemispheric wedge $P_W$, the KMS state at $\beta = 2\pi$
on the wedge sub-algebra, and the unified surface-gravity normalisation.
Section~\ref{sec:geometric_construction} constructs the posited geometric
BW-$\alpha$ generator $K_\alpha = J_{\mathrm{polar}}$ (its
identification as a modular Hamiltonian is open, O0) and
proves the automatic period closure Theorem~\ref{thm:bw_alpha}.
Section~\ref{sec:tt_construction} constructs the Tomita--Takesaki
BW-$\gamma$ modular Hamiltonian $K_{\TT} = -\log \Delta$ of the Gibbs
state built from that posited generator, via polar
decomposition on the GNS Hilbert--Schmidt space, distinguishes
$J_{\TT}$ ($J^2 = +I$) from the Connes $J_{\mathrm{GV}}$ ($J^2 = -I$),
and proves Theorem~\ref{thm:bw_gamma} (the same closure, read through
the Gibbs state built from the posited generator).
Section~\ref{sec:unified_closure} proves
Theorem~\ref{thm:unified} (the flow conjugacy $\sigma_t^{\TT} =
\sigma_{-t}^{\alpha}$, Tomita's theorem for a Gibbs state) and
discusses the local Hamiltonian $H_{\mathrm{local}} = K_\alpha^W /
\beta$ as a posited choice.
Section~\ref{sec:four_witness_collapse} records that the HH, Sewell,
BW, and Unruh instantiations coincide by construction. Section~\ref{sec:related_work} places the result in the
modular-theory and NCG literature. Section~\ref{sec:conclusion}
summarises and names open questions.
Section~\ref{sec:lorentzian_closure} carries the construction onto a
signature-$(3, 1)$ Krein carrier at finite cutoff (the successor of the
former Paper~43).
Appendix~\ref{app:computational_verification} documents the
computational verification protocol and the supporting module
\texttt{geovac/modular\_hamiltonian.py}.
Appendix~\ref{app:framework_crossref} cross-references the result to
the present author's broader framework
(Paper~38~\cite{paper38} for state-space Gromov--Hausdorff convergence;\
Paper~40~\cite{paper40_unified} for the universal $4/\pi$ rate
constant;\ Paper~32~\S~VIII for the master Mellin engine
case-exhaustion theorem and the four-witness theorem's M1
sub-mechanism reading).

% =====================================================================
\section{The truncated Camporesi--Higuchi spectral triple}\label{sec:setup}
% =====================================================================

We work in the framework of metric spectral triples and their
spectral truncations established by
Connes~\cite{connes1995, connes_vs2021} and developed in the
GH-convergence context by
Latr\'emoli\`ere~\cite{latremoliere2018, latremoliere_metric_st_2017},
Hekkelman~\cite{hekkelman2022},
Leimbach--van~Suijlekom~\cite{leimbach_vs2024}, and the present
author's Paper~38~\cite{paper38} for the round $\sthree = \SU(2)$
case.

\subsection{The continuum Camporesi--Higuchi triple}
\label{sec:ch_triple_recap}

Identify the unit round $\sthree$ with the Lie group $\SU(2)$ via
the standard double cover, equipped with the bi-invariant metric of
total volume $\Vol(\sthree) = 2\pi^2$. Let $\Sigma$ denote the spinor
bundle and $L^2(\sthree, \Sigma)$ the spinor Hilbert space. The
\emph{Camporesi--Higuchi Dirac operator}
$\DCH: \mathrm{Dom}(\DCH) \subset L^2(\sthree, \Sigma) \to L^2(\sthree,
\Sigma)$ acts diagonally on the spinor harmonic basis
$\{\psi_{n, l, m_j, \chi}\}$ via
\begin{equation}\label{eq:CH_spectrum}
   \DCH \, \psi_{n, l, m_j, \chi}
   \;=\; \chi\,(n + \tfrac{1}{2})\, \psi_{n, l, m_j, \chi},
\end{equation}
with $n \in \N_{\ge 1}$, $l \in \{0, 1, \ldots, n-1\}$,
$j = l + 1/2$, $m_j \in \{-j, -j+1, \ldots, j\}$, and chirality
$\chi \in \{+1, -1\}$~\cite{camporesi_higuchi1996}. The degeneracy of
the level $n$ on each chirality sector is $n(n+1)$;\ the total
degeneracy including both chiralities is $2n(n+1)$.

We form the spectral triple
\begin{equation}\label{eq:T_S3}
   \Tcal_{\sthree}
   \;=\; \bigl(\, C^\infty(\sthree),\;
                 L^2(\sthree, \Sigma),\;
                 \DCH \,\bigr).
\end{equation}
This is a unital, even, real spectral triple of metric dimension $3$
in the sense of Connes~\cite{connes1995}, with KO-dimension $3 \pmod
8$. The chirality grading $\gamma_5 = \mathrm{diag}(+1, -1)$ on
$\Sigma_+ \oplus \Sigma_-$, the Connes real structure
$J_{\mathrm{GV}}$ (charge conjugation, $J_{\mathrm{GV}}^2 = -I$), and
the standard Lipschitz seminorm
$\opnorm{[\DCH, \Mat_f]}$ on the algebra $C^\infty(\sthree)$ all enter
as expected. We use the labelling $n_{\mathrm{Fock}} = n_{\mathrm{CH}}
+ 1$ throughout, so the lowest non-trivial level is $n = 1$ with
$|\lambda_1| = 3/2$.

\subsection{Truncations and the operator system}
\label{sec:truncations_recap}

Let $P_{\nmax}: L^2(\sthree, \Sigma) \to \Hilb_{\nmax}$ denote the
orthogonal projection onto the finite-dimensional subspace spanned by
the spinor harmonics with $1 \le n \le \nmax$. The truncated Hilbert
space has dimension
\begin{equation}\label{eq:dim_H_nmax}
   \dim \Hilb_{\nmax}
   \;=\; 2 \sum_{n=1}^{\nmax} n(n+1)
   \;=\; \tfrac{2}{3}\, \nmax(\nmax+1)(\nmax+2).
\end{equation}
Numerically:\ $\dim \Hilb_2 = 16$, $\dim \Hilb_3 = 40$,
$\dim \Hilb_4 = 80$, $\dim \Hilb_5 = 140$. The $\nmax = 3$ case is
numerically distinguished by the coincidence $\dim \Hilb_3 = 40 =
g_3^{\mathrm{Dirac}} = \Delta^{-1}$ --- different objects:\
$\dim \Hilb_3 = 4 + 12 + 24$ is the cumulative node count over the
Camporesi--Higuchi levels $k \le 2$, while $g_3^{\mathrm{Dirac}} =
2(3+1)(3+2) = 40$ is the both-sign Dirac degeneracy of the single
level $k = 3$ ($|\lambda| = 9/2$); they agree only at $\nmax = 3$ ---
where $\Delta^{-1}$ is the boundary-mass invariant of the present
author's Paper~2~\cite{paper2} fine-structure-constant observation.

The truncated Dirac operator is $D_{\nmax} = P_{\nmax} \DCH P_{\nmax}$,
which (since $\DCH$ is diagonal in the spinor basis) acts as
$\chi\,(n + 1/2)$ on each kept basis vector. Following
Connes--van~Suijlekom~\cite{connes_vs2021}, the \emph{truncated
operator system} is
\begin{equation}\label{eq:O_def}
   \Op_{\nmax}
   \;=\; P_{\nmax}\, C^\infty(\sthree)\, P_{\nmax}
   \;\subset\; B(\Hilb_{\nmax}),
\end{equation}
viewed as a self-adjoint subspace of the matrix algebra
$\Mat_{\dim \Hilb_{\nmax}}(\C)$. As emphasised in
\cite[\S 1]{connes_vs2021}, this is \emph{not} a $C^*$-subalgebra of
$\Mat_{\dim \Hilb_{\nmax}}(\C)$:\ the full $C^*$-algebra it generates
is the full matrix algebra (the ``$C^*$-envelope''), but the operator
system $\Op_{\nmax}$ itself encodes finer structure, including the
algebra-side analog of the Lipschitz seminorm and the propagation
invariant. The truncated metric spectral triple is then
$\Tcal_{\nmax} = (\Op_{\nmax}, \Hilb_{\nmax}, D_{\nmax})$.

Paper~32~\S~III of the present author's framework documentation
verifies that the propagation invariant
$\mathrm{prop}(\Op_{\nmax}) = 2$ at $\nmax \in \{2, 3, 4\}$, matching
Connes--van~Suijlekom~\cite[\S 4]{connes_vs2021} for the
Toeplitz operator system on $S^1$ verbatim. This is the structural
match that establishes $\Op_{\nmax}$ as the right operator-system
host for spectral-triple-internal modular constructions.

\subsection{State-space GH convergence to the continuum}
\label{sec:propinquity_recap}

The present author's Paper~38~\cite{paper38} proves that the
truncated triples $\Tcal_{\nmax}$ converge to $\Tcal_{\sthree}$ in
van~Suijlekom's state-space Gromov--Hausdorff
distance with explicit rate
\begin{equation}\label{eq:paper38_main}
   \Lambda(\Tcal_{\nmax}, \Tcal_{\sthree})
   \;\le\; C_3 \cdot \gamma_{\nmax},
   \qquad
   \gamma_{\nmax} \;=\; \frac{4 \log \nmax}{\pi \, \nmax}
                       \;+\; O(1/\nmax),
\end{equation}
with Lipschitz comparison constant $C_3 = 1$ and asymptotic rate
constant $4/\pi = \Vol(S^2)/\pi^2$ (in the dual-Coxeter normalisation
of the $\SU(2)$ metric;\ $2/\pi$ on the unit round $\sthree$, the
constant scaling with the metric). It carries the $k = 0$ volume
content of the M1 sub-mechanism of the master Mellin engine.  It is
\emph{not} the Hopf fibration's base-to-total ratio and not an
$\SU(2)/U(1)$ Haar quotient:\ the Hopf map is a Riemannian submersion
$S^3(r) \to S^2(r/2)$, whose ratio is neither $4/\pi$ nor $2/\pi$;\ the
constant is a quotient of \emph{unit-radius} sphere volumes,
$\Vol(S^2)/\Vol(\sthree)$.  (The historical label ``Hopf-base measure''
remains in use for the $k = 0$ Mellin slot itself and for Paper~25's
$\Vol(S^2)/4 = \pi$;\ what was retired on 2026-09-03 is reading the
rate constant as a base-to-total ratio.  Paper~38, volume-ratio note.)
The M1 signature is on $\sthree$ (Paper~32~\S~VIII case-exhaustion theorem;
generalised to compact connected semisimple Lie groups in
Paper~40~\cite{paper40_unified}; a circle factor carries the circle constant). The lemmas L1' / L2 / L3 / L4
of Paper~38 — chirality-doubled operator system substrate, central
spectral Fej\'er kernel on $\SU(2)$, Lipschitz comparison bound,
Berezin reconstruction — supply the analytic infrastructure that the
present paper builds on.  (Paper~38's fifth lemma, the
Latr\'emoli\`ere propinquity assembly, was \emph{withdrawn} on
2026-09-03 when its height leg was refuted;\ nothing here depends on
it, and Paper~38's convergence theorem is carried by the
lifted-state construction instead.)

The present paper does not require the GH-convergence machinery in any
load-bearing way:\ the period closure proved below is a
finite-cutoff bit-exact identity, not a limit statement. The
state-space Gromov--Hausdorff convergence~\eqref{eq:paper38_main}
enters only as a
\emph{consistency} sanity check (Section~\ref{sec:tt_construction});
it does not link the modular period to the metric-side rate constant.

% =====================================================================
\section{Lorentzian-readiness audit and Riemannian scope}
\label{sec:lorentzian_readiness}
% =====================================================================

The setting of the present paper is signature $(3, 0)$:\ the round
$\sthree$ Camporesi--Higuchi spectral triple is Riemannian, with
KO-dimension $3 \pmod 8$. The four-witness Wick-rotation theorem
inhabits a hybrid status:\ each of the four physical witnesses
(Hartle--Hawking, Sewell, Bisognano--Wichmann, Unruh) is naturally a
\emph{Lorentzian} statement at signature $(3, 1)$ in the Wightman
axiomatic framework, but the shared $\beta = 2\pi / \kappa_g$ structural
identity that unifies the four faces lives entirely on the Euclidean
side via Wick rotation. The Wick-rotation chain itself
[Hartle--Hawking~1976~$\to$~Sewell~1982~$\to$~Bisognano--Wichmann~1976~$\to$~Unruh~1976]
operates at the metric-functional level:\ it converts a Euclidean
heat-kernel-style spectral computation into its Lorentzian
modular-flow interpretation without re-deriving the underlying
spectral object in Krein signature.

\subsection{The 28-projection Lorentzian transfer audit}
\label{sec:l0_audit}

The present author's Sprint L0 audit
(\texttt{debug/lorentzian\_l0\_audit\_memo.md}, Paper~31~\S~8,
Paper~34~\S~V.E) classifies the 28 projections of the framework's
projection taxonomy into four buckets:
\begin{enumerate}\setlength{\itemsep}{2pt}
\item[(T)] \emph{TRANSFERS\_FREELY:}\ signature-blind, lives on the
variable/dimension axes of the projection taxonomy only, with no
Riemannian volume form, no Mellin transform, and no Latr\'emoli\`ere
propinquity. Holds at $(3, 0)$ and $(3, 1)$ identically.
\item[(W)] \emph{WICK\_MAP\_FREE:}\ reduces to the M1 sub-mechanism
of the master Mellin engine via the four-witness Wick-rotation theorem
at the metric-functional level.
\item[(E)] \emph{EUCLIDEAN\_SPECIFIC:}\ requires either (i)~a
Krein-space spectral triple at $(3, 1)$ per
Bizi--Brouder--Besnard~2018~\cite{bizi_brouder_besnard2018}, (ii)~a
Lorentzian-native heat-trace formalism, or (iii)~a Lorentzian
analog of the Latr\'emoli\`ere propinquity (which does not yet exist
in the published literature).
\item[(M)] \emph{MIXED:}\ parts transfer freely, parts require
extension.
\end{enumerate}

The headline distribution at May 2026 is
$17 / 4 / 5 / 2 = 61\% / 14\% / 18\% / 7\%$: three-quarters of the
projection dictionary inherits a Lorentzian reading at no further
structural cost at the metric-functional level. The Wick-rotation
projection itself sits in the (W) bucket and carries the M1
transcendental signature $2\pi \cdot \Q$ via $\Vol(S^1)$. The
five EUCLIDEAN\_SPECIFIC projections (Fock conformal, Hopf bundle,
stereographic, spectral action, Camporesi--Higuchi spinor lift) all
require a full Krein-space lift to extend cleanly to signature
$(3, 1)$.

Sections~\ref{sec:setup}--\ref{sec:conclusion} do not address signature
$(3, 1)$ (the Krein carrier of \S\ref{sec:lorentzian_closure} does, at
finite cutoff). They operate
entirely on the Riemannian Camporesi--Higuchi triple at signature
$(3, 0)$ and record the modular-period identity inside that framework.
The Sprint L0 audit identified an operator-system test of the
Wick-rotation projection at signature $(3, 0)$ --- a modular
Hamiltonian computed from a wedge-restricted Camporesi--Higuchi state
--- as the principal Sprint L1 falsifier ahead of the $(3, 1)$
extension (itself beyond sprint scale).  The present paper does not
run that test:\ its generator is posited, so its period closure is
automatic (withdrawn as a closure of the falsifier, 2026-09-28).

\subsection{Footnote on the Nieuviarts twist morphism}
\label{sec:nieuviarts_footnote}

A natural shortcut to the signature change would be the
twisted-spectral-triple morphism of
Nieuviarts~\cite{nieuviarts2025a, nieuviarts2024, nieuviarts2025b_proceedings},
which constructs a pseudo-Riemannian spectral triple from an
almost-commutative twisted Riemannian spectral triple, presented as
an alternative to Wick rotation. The construction is mechanistically
appealing but does \emph{not} apply to the present $\sthree$ case:\
Definition~2.2 of~\cite{nieuviarts2025a} requires a $\Z_2$ grading
$\Gamma$ that commutes with the algebra and anti-commutes with the
Dirac, which exists only for \emph{even} KO-dimension spectral
triples. The $\sthree$ Camporesi--Higuchi triple has KO-dimension
$3 \pmod 8$ (odd); the restriction is structural rather than
editorial (the author of~\cite{nieuviarts2025a} explicitly notes:
``the fact that no such procedure for twisting odd-dimensional
manifold's spectral triple has been found will be one of the
justifications to focus on the study of even-dimensional manifolds'').
The Lorentzian-extension path for the present construction therefore
runs through the direct Bizi--Brouder--Besnard~2018 Krein-space lift,
not through a twist morphism;\ this is the named Sprint L2 follow-up
of the present author's framework documentation. A dedicated NO-GO
analysis of the Nieuviarts route on $\sthree = \SU(2)$ is recorded in
\texttt{debug/nieuviarts\_scoping\_memo.md} of the present author's
internal materials.

% =====================================================================
\section{The Bisognano--Wichmann modular structure at finite cutoff}
\label{sec:bw_at_finite_cutoff}
% =====================================================================

We set up the operator-system-level analog of the
Bisognano--Wichmann construction on the truncated Camporesi--Higuchi
spectral triple. Three ingredients are required:\ (i)~the hemispheric
wedge projector $P_W$, (ii)~the wedge-restricted Gibbs state
$\omega_W$ at inverse temperature $\beta$, (iii)~the local Hamiltonian
$H_{\mathrm{local}}$ generating the wedge-modular evolution.

\subsection{The hemispheric wedge projector}\label{sec:hemispheric_wedge}

\begin{definition}[Hemispheric wedge on $\Hilb_{\nmax}$]
\label{def:wedge}
The \emph{Hopf-axis equatorial reflection} on $\Hilb_{\nmax}$ is the
involution
\begin{equation}\label{eq:R_polar_def}
   R_{\mathrm{polar}}: \psi_{n, l, m_j, \chi}
   \;\longmapsto\; \psi_{n, l, -m_j, \chi},
\end{equation}
acting as $m_j \mapsto -m_j$ on the spinor basis. The
\emph{hemispheric wedge projector} is
\begin{equation}\label{eq:P_W_def}
   P_W \;:=\; \tfrac{1}{2}(I + R_{\mathrm{polar}})
   \;\in\; B(\Hilb_{\nmax}).
\end{equation}
\end{definition}

\begin{proposition}[Wedge projector properties]\label{prop:wedge_properties}
The hemispheric wedge projector $P_W$ satisfies:
\begin{enumerate}\setlength{\itemsep}{2pt}
\item[(a)] $P_W^2 = P_W$ (idempotent);
\item[(b)] $P_W^* = P_W$ (Hermitian);
\item[(c)] $\Tr(P_W) = \dim \Hilb_{\nmax} / 2$ (half-dimensional
projection).
\end{enumerate}
Equivalently, $\Hilb_{\nmax} = \Hilb_W \oplus \Hilb_{W'}$ where
$\Hilb_W = P_W \Hilb_{\nmax}$ has dimension $\dim_W = \dim \Hilb_{\nmax}/2$,
i.e., $\dim_W = 8, 20, 40, 70$ at $\nmax = 2, 3, 4, 5$.
\end{proposition}

\begin{proof}
(a) and (b) follow from $R_{\mathrm{polar}}^2 = I$ and
$R_{\mathrm{polar}}^* = R_{\mathrm{polar}}$, the latter holding because
$R_{\mathrm{polar}}$ is a permutation matrix. For (c):\ on the
full-Dirac spinor basis the values $\mathrm{two}\_m_j$ are
\emph{odd integers}, so the involution $m_j \mapsto -m_j$ has no fixed
points, i.e., $R_{\mathrm{polar}}$ has spectrum $\{+1, -1\}$ each with
multiplicity $\dim \Hilb_{\nmax}/2$.
\end{proof}

The hemispheric wedge $\Hilb_W$ is the operator-system analog of the
half-$\sthree$ algebra $C^\infty(H_+)$ where $H_+ = \{\omega \in
\sthree : \omega \cdot \hat{n} > 0\}$ is the open hemisphere of
$\sthree$ aligned with a chosen polar axis $\hat{n}$. The choice of
axis is fixed canonically by alignment with the Hopf-base direction
(see~\cite[\S V.3]{paper25_hopf_gauge_structure} for the role of the
Hopf-base axis in the present author's framework documentation, and
Section~\ref{sec:tt_construction} below for the role of the integer
spectrum of $\mathrm{two}\_m_j$ on this axis).

\begin{remark}[On the choice of wedge]\label{rem:wedge_choice}
Three candidates for the wedge sub-system were named in the architectural
scoping (\texttt{debug/l1\_modular\_hamiltonian\_architecture\_memo.md}
\S 3):\ (W1) the hemispheric wedge of Definition~\ref{def:wedge};
(W2) the $\kappa$-preserving block of Paper~29's Dirac graph Rule A;
(W3) the pendant-vertex-free subgraph of the scalar Fock graph.
Candidate W1 is the cleanest Riemannian-side analog of the Wightman
Rindler wedge (a region bounded by a horizon of codimension 1 in the
underlying manifold), and it carries a natural Killing-preserving
structure (the rotation around the Hopf-base axis preserves the
wedge boundary). Candidates W2 and W3 are graph-topological
projections rather than spatial-region restrictions, and their modular
structure either degenerates (W2) or has no clear ``localised in a
half-region'' interpretation (W3). The present paper proceeds with W1
as the canonical choice. The dependence of the closure on the wedge
choice is part of the open scope (Section~\ref{sec:conclusion}).
\end{remark}

\subsection{The wedge-restricted Gibbs state}\label{sec:kms_state}

For a self-adjoint operator $H$ on $\Hilb_W$ and a positive inverse
temperature $\beta > 0$, the Gibbs state
$\omega_{\beta, H}: B(\Hilb_W) \to \C$ is
\begin{equation}\label{eq:gibbs_state_def}
   \omega_{\beta, H}(O)
   \;=\; \frac{\Tr_{\Hilb_W}\bigl(e^{-\beta H}\, O\bigr)}
              {\Tr_{\Hilb_W}\bigl(e^{-\beta H}\bigr)},
   \qquad O \in B(\Hilb_W).
\end{equation}
Restricted to the wedge sub-operator-system
$\Op_{W, \nmax} := P_W \Op_{\nmax} P_W$, this is the natural
Hilbert-space-internal analog of the Wightman Rindler-wedge KMS state.

The choice of $H$ is the load-bearing structural input, and it is a
\emph{posit}:\ nothing below derives it from the Camporesi--Higuchi
triple. We make the choice suggested by the form of the Wightman BW
vacuum (Section~\ref{sec:tt_construction} below):
\begin{equation}\label{eq:H_local_def}
   H_{\mathrm{local}} \;:=\; K_\alpha^W \,/\, \beta,
\end{equation}
where $K_\alpha^W$ is the geometric BW-$\alpha$ generator of
Definition~\ref{def:K_alpha_W} below. With this choice, $\beta
H_{\mathrm{local}} = K_\alpha^W$ is itself $\beta$-independent at the
algebra-action level, so the wedge KMS state $\rho_W = e^{-K_\alpha^W}/Z$
is the same for every $\beta$. This is a deliberate construction-level
choice;\ an alternative choice $H_{\mathrm{local}} = D_W$ (the
wedge-restricted Camporesi--Higuchi Dirac) would lift the $\beta = 2\pi$
period to the spinor double-cover $-I$ rather than the identity
operator $+I$ of Section~\ref{sec:geometric_construction}
(the conjugation flow still closes, but at the double-cover sign).
We return to this point in Section~\ref{sec:unified_closure}.

\subsection{Four-witness unification and unit normalisation}
\label{sec:four_witness_unit}

The four physical witnesses (Hartle--Hawking, Sewell,
Bisognano--Wichmann, Unruh) share the structural identity
\begin{equation}\label{eq:beta_kappa_g}
   \beta_W \;=\; 2\pi \,/\, \kappa_g(W),
\end{equation}
where $\kappa_g(W)$ is the surface gravity of the bifurcate Killing
horizon bounding the causal region $W$:\
$\kappa_g^{\HH} = 1/(4M)$ for the Schwarzschild exterior of mass $M$,
giving $\beta^{\HH} = 8\pi M$;\
$\kappa_g^{\Sew} = \kappa_g^{\HH}$ for the abstract KMS-modular
reading of the Hartle--Hawking state;\
$\kappa_g^{\BW} = 1$ in natural rapidity units for the Rindler wedge,
giving $\beta^{\BW} = 2\pi$;\
$\kappa_g^{\Unruh} = a$ for the Rindler wedge in proper-acceleration
units, giving $\beta^{\Unruh} = 2\pi / a$.

\paragraph{Unit normalisation.}
We adopt the BW rapidity-units convention $\kappa_g = 1$ (denoted
\emph{U-1} in the architectural scoping) as the BW witness.
The four physical witnesses are then instantiated as parameterised
BW with witness-specific $\kappa_g$ values: $\BW \to \kappa_g = 1$,
$\HH \to \kappa_g = 1/(4M)$, $\Sew \to \kappa_g = 1/(4M)$,
$\Unruh \to \kappa_g = a$. At the operator level every instantiation
uses the \emph{same} generator $K_\alpha^W$;\ $\kappa_g$ enters only as
a label (through $\beta = 2\pi/\kappa_g$ and $H_{\mathrm{local}} =
K_\alpha^W/\beta$), and cancels in $\beta H_{\mathrm{local}}$.  The
period closure $\sigma_{2\pi}(O) = O$ is therefore
$\kappa_g$-independent by construction
(Sections~\ref{sec:geometric_construction}--\ref{sec:four_witness_collapse}),
and the agreement of the four witnesses at finite cutoff is a
consistency check of the bookkeeping, not a finite-cutoff theorem about
them.

% =====================================================================
\section{The geometric construction (BW-$\alpha$)}
\label{sec:geometric_construction}
% =====================================================================

The first construction on $\Tcal_{\nmax}$ is the geometric
BW-$\alpha$ generator, obtained by \emph{positing} that the role of the
wedge-preserving boost is played by the spatial rotation generator
around the Hopf-base axis.  On the Riemannian $\sthree$ there is no
boost;\ the rotation generator is a stand-in, not an identification.

\subsection{The geometric generator}\label{sec:K_alpha_W}

\begin{definition}[BW-$\alpha$ geometric generator on the wedge]
\label{def:K_alpha_W}
The \emph{BW-$\alpha$ geometric generator} on the wedge sub-Hilbert
space $\Hilb_W = P_W \Hilb_{\nmax}$ is the operator
\begin{equation}\label{eq:K_alpha_W_def}
   K_\alpha^W \;:=\; \mathrm{diag}\bigl(\mathrm{two}\_m_j\bigr)
   \quad \text{on} \quad \Hilb_W,
\end{equation}
acting on each wedge basis state $\psi_{n, l, |m_j|, \chi}$ (with
$m_j$ chosen positive on the wedge $P_W$) by multiplication by
$\mathrm{two}\_m_j$, the doubled $m$-quantum-number.
\end{definition}

\begin{remark}[Why $\mathrm{two}\_m_j$ rather than $m_j$]
\label{rem:two_m_j}
On the full-Dirac spinor basis $\{\psi_{n, l, m_j, \chi}\}$, the
quantum number $m_j$ takes half-integer values
$\{-l - 1/2, -l + 1/2, \ldots, l + 1/2\}$. To obtain an integer
spectrum --- the load-bearing property for the bit-exact period
closure of Theorem~\ref{thm:bw_alpha} --- we work with the doubled
generator $\mathrm{two}\_m_j = 2 m_j$, taking odd integer values.
This is the standard half-integer-spin doubling convention from
representation theory of $\SU(2)$. The factor of $2$ does not affect
the modular flow at the period (modular flow is invariant under
overall rescaling of $K$ by integer factors), but it makes the
integer-spectrum property manifest.
\end{remark}

\begin{proposition}[Integer spectrum]\label{prop:integer_spectrum}
The spectrum of $K_\alpha^W$ on $\Hilb_W$ is contained in the odd
integers $\{-(2l+1), -(2l-1), \ldots, +(2l-1), +(2l+1)\}$ over
$l \in \{0, 1, \ldots, \nmax - 1\}$. In particular,
$\mathrm{Spec}(K_\alpha^W) \subset \Z$.
\end{proposition}

\begin{proof}
On the wedge $P_W$, $m_j$ takes values $\{1/2, 3/2, \ldots, l + 1/2\}$
(positive half-integers up to $l + 1/2$). The doubled generator
$\mathrm{two}\_m_j$ takes values $\{1, 3, \ldots, 2l + 1\}$, all odd
positive integers. Including the wedge-anti-wedge structure via the
full-Dirac basis would give the full range $\{-(2l+1), \ldots, +(2l+1)\}$
in odd integer steps; on the wedge itself we have the positive
sub-range.
\end{proof}

\subsection{The geometric modular flow}\label{sec:alpha_flow}

The geometric BW-$\alpha$ flow (a modular flow only if O0 identifies
$K_\alpha^W$ as a modular Hamiltonian) on the wedge sub-operator-system
$\Op_{W, \nmax}$ is the inner-derivation flow
\begin{equation}\label{eq:alpha_flow}
   \sigma_t^{\alpha}(O)
   \;:=\; e^{i t K_\alpha^W}\, O\, e^{-i t K_\alpha^W},
   \qquad O \in \Op_{W, \nmax},\; t \in \R.
\end{equation}
This is well-defined as a strongly continuous one-parameter group of
unital $*$-automorphisms on $B(\Hilb_W)$. Its restriction to
$\Op_{W, \nmax}$ may leave the operator system (since $\Op_{W, \nmax}$
is not multiplicatively closed); we return to this point in
Section~\ref{sec:os_leakage}.

\begin{theorem}[BW-$\alpha$ period closure at finite cutoff, {[SYMBOLIC]}]
\label{thm:bw_alpha}
For every $\nmax \ge 1$ and every $O \in B(\Hilb_W)$, the BW-$\alpha$
modular flow~\eqref{eq:alpha_flow} satisfies the period closure
\begin{equation}\label{eq:bw_alpha_closure}
   \sigma_{2\pi}^{\alpha}(O) \;=\; O
\end{equation}
identically, i.e., as a bit-exact operator identity.
\end{theorem}

\begin{proof}
By Proposition~\ref{prop:integer_spectrum}, $K_\alpha^W$ is diagonal
with integer eigenvalues $\{n_1, n_2, \ldots, n_{\dim_W}\}$ in its
eigenbasis. The unitary $e^{i \cdot 2\pi \cdot K_\alpha^W}$ is then
diagonal with eigenvalues $e^{i \cdot 2\pi \cdot n_j} = 1$ for every
$j$ (since each $n_j$ is an integer):\
$e^{i \cdot 2\pi \cdot K_\alpha^W} = I_{\dim_W}$ identically. Hence
$\sigma_{2\pi}^{\alpha}(O) = I \cdot O \cdot I = O$ for every $O$.
\end{proof}

\begin{remark}[What Theorem~\ref{thm:bw_alpha} does not test]
\label{rem:bw_alpha_automatic}
The proof uses only that $K_\alpha^W$ is Hermitian with integer
spectrum;\ it uses nothing about the wedge, the Camporesi--Higuchi
Dirac, or any state.  The same closure holds for \emph{every}
integer-spectrum generator and \emph{every} operator $O$, so it cannot
distinguish a Bisognano--Wichmann modular flow from any other
$2\pi$-periodic unitary group.  (Concretely:\ replacing $K_\alpha^W$ in
the supporting module by an unrelated integer-spectrum matrix leaves
green all eighteen period-closure and cross-witness tests of
\texttt{tests/test\_modular\_hamiltonian.py} exercised in that fire
test, 2026-09-28.)  The theorem is a correct elementary fact;\ it is not evidence
that $K_\alpha^W$ is the modular Hamiltonian of a physically
distinguished wedge state.
\end{remark}

\paragraph{Numerical verification.}
Theorem~\ref{thm:bw_alpha} is an exact algebraic identity at the level
of the symbolic construction. In a floating-point computation, the
matrix exponential $e^{i \cdot 2\pi \cdot K_\alpha^W}$ is computed by
diagonalisation and then exponentiation of the eigenvalues, so the
identity holds to within the floating-point precision of the
diagonalisation step. We have computed the maximum periodicity
residual
\begin{equation}\label{eq:periodicity_residual_def}
   r_W^{\alpha}(O) \;:=\; \Frob{\sigma_{2\pi}^{\alpha}(O) - O}
\end{equation}
for the first five non-identity multipliers $O \in \Op_{\nmax}$
(restricted to the wedge), at $\nmax \in \{2, 3, 4, 5\}$. The
results are recorded in Table~\ref{tab:bw_alpha_residuals}.

\begin{table}[ht]
\centering
\caption{Maximum BW-$\alpha$ periodicity residual
$\max_O r_W^{\alpha}(O)$ at finite cutoff. Computed in
double-precision floating point on the full-Dirac basis;
residuals scale as
$O(\sqrt{\dim \Hilb_{\nmax}} \cdot \varepsilon_{\mathrm{machine}})$,
i.e., pure round-off accumulation with no structural drift. The
verdict column reproduces the supporting module's label;\ it records
that the automatic closure was reproduced, not an identification
(Remark~\ref{rem:bw_alpha_automatic}). Source
data:\ \texttt{debug/data/l1\_modular\_hamiltonian\_results.json}.}
\label{tab:bw_alpha_residuals}
\begin{tabular}{rrrcc}
\toprule
$\nmax$ & $\dim \Hilb_{\nmax}$ & $\dim_W$ &
$\max_O r_W^{\alpha}(O)$ &
verdict\\
\midrule
2 & 16  & 8  & $1.80 \times 10^{-16}$ & STRONG\_IDENTIFICATION \\
3 & 40  & 20 & $4.76 \times 10^{-16}$ & STRONG\_IDENTIFICATION \\
4 & 80  & 40 & $9.33 \times 10^{-16}$ & STRONG\_IDENTIFICATION \\
5 & 140 & 70 & $1.59 \times 10^{-15}$ & STRONG\_IDENTIFICATION \\
\bottomrule
\end{tabular}
\end{table}

The $\nmax = 3$ row is the one at which $\dim \Hilb_3 = 40$
numerically coincides with $g_3^{\mathrm{Dirac}} = \Delta^{-1} = 40$
(different objects that agree only at $\nmax = 3$, as noted above),
the selection-principle cutoff of the present author's
Paper~2~\cite{paper2} fine-structure-constant observation. Nothing
about the closure is special at this cutoff:\ it holds at every
$\nmax$ for the integer-spectrum reason above.

\subsection{Operator-system non-leakage}\label{sec:os_leakage}

A separate concern flagged in the architectural scoping
(\texttt{debug/l1\_modular\_hamiltonian\_architecture\_memo.md}
\S 10a) is that the truncated operator system $\Op_{\nmax}$ is
$*$-closed but not multiplicatively closed:\ the modular flow
$\sigma_t^{\alpha}$ acts on all of $B(\Hilb_W)$, and its restriction
to $\Op_{W, \nmax}$ might leak out. The integer-spectrum closure
Theorem~\ref{thm:bw_alpha} resolves this without further work:\
$\sigma_{2\pi}^{\alpha}(O) = O$ for every $O$, so the flow at the
period $t = 2\pi$ leaves every element of $\Op_{W, \nmax}$ inside
$\Op_{W, \nmax}$ trivially. At intermediate $t \in (0, 2\pi)$ the flow
may leave the operator system, but the period closure (the modular
identity we test) is unaffected.

% =====================================================================
\section{The Tomita--Takesaki construction (BW-$\gamma$)}
\label{sec:tt_construction}
% =====================================================================

The second construction on $\Tcal_{\nmax}$ is the Tomita--Takesaki
BW-$\gamma$ reading, obtained by polar decomposition of the modular $S$
operator on the GNS Hilbert--Schmidt space of the wedge-restricted
Gibbs state.  Because that state is built from the posited generator,
$\rho_W = e^{-K_\alpha^W}/Z$, this construction returns $K_\alpha^W$
(Tomita's theorem for a Gibbs state) and is \emph{not} an independent
check on it;\ its value here is that it exhibits the GNS/Tomita
machinery explicitly at finite cutoff, ready for the [OPEN] case in
which the state is not built from $K$.

\subsection{The GNS Hilbert--Schmidt space}\label{sec:gns}

Let $\rho := e^{-K_\alpha^W}/Z$ be the wedge density matrix
(Section~\ref{sec:kms_state} with $H_{\mathrm{local}} = K_\alpha^W /
\beta$, so $\beta H_{\mathrm{local}} = K_\alpha^W$ is
$\beta$-independent). For the von Neumann algebra
$\Acal_W := B(\Hilb_W) \cong M_{\dim_W}(\C)$ (Type I$_{\dim_W}$
factor) and the faithful normal state
$\omega(O) := \Tr_{\Hilb_W}(\rho \, O)$, the GNS triple
$(\pi_\omega, \Hilb_{\GNS}, \Omega)$ is realised on the
Hilbert--Schmidt space
\begin{equation}\label{eq:H_GNS_def}
   \Hilb_{\GNS}
   \;=\; M_{\dim_W}(\C)
   \;\cong\; \C^{\dim_W^2},
   \qquad
   \langle X, Y \rangle_{\GNS} \;:=\; \Tr(X^* Y),
\end{equation}
with cyclic vector $\Omega := \rho^{1/2}$ (the canonical purification
of $\rho$). The left action of $a \in \Acal_W$ is $L_a(X) = a X$, the
right action of $a$ from the commutant $\Acal_W^{\mathrm{op}}$ is
$R_a(X) = X a$. Numerically, $\dim \Hilb_{\GNS} = \dim_W^2 = 64, 400,
1600, 4900$ at $\nmax = 2, 3, 4, 5$.

\subsection{The modular $S$ operator and polar decomposition}
\label{sec:S_polar}

The Tomita modular operator $S: \pi_\omega(a)\Omega \mapsto
\pi_\omega(a^*) \Omega$ is the antilinear closure of the conjugate-linear
involution $a \Omega \mapsto a^* \Omega$ on the cyclic-separating
orbit. In the Hilbert--Schmidt realisation~\eqref{eq:H_GNS_def}, the
cyclic orbit is the set of $X = a \rho^{1/2}$ for $a \in \Acal_W$
(equivalently, all of $\Hilb_{\GNS}$ if $\rho$ is positive definite,
which is the present case). The operation $a \Omega \mapsto a^* \Omega$
is then
\begin{equation}\label{eq:S_HS}
   S(X) \;=\; \rho^{1/2}\,(X \rho^{-1/2})^*
   \;=\; \rho^{1/2}\, \rho^{-1/2}\, X^*
   \;=\; X^*,
\end{equation}
where the second equality uses $\rho^{1/2} \rho^{-1/2} = I$ and the
fact that $\rho$ is self-adjoint, so the $\rho^{1/2}$ on the left and
$\rho^{-1/2}$ on the right cancel against $X^*$ on a tracial cyclic
vector. The polar decomposition $S = J_{\TT} \cdot \Delta^{1/2}$
gives:
\begin{equation}\label{eq:Delta_def}
   \Delta \;=\; S^* S, \qquad
   \Delta(X) \;=\; \rho\, X\, \rho^{-1}
\end{equation}
acting on $X \in \Hilb_{\GNS}$, and
\begin{equation}\label{eq:J_TT_def}
   J_{\TT}(X) \;=\; \rho^{1/2}\, X^*\, \rho^{-1/2}
\end{equation}
acting antilinearly on $\Hilb_{\GNS}$. In the column-stacked matrix
form $\mathrm{vec}(X) \in \C^{\dim_W^2}$, the modular operator
$\Delta$ is
\begin{equation}\label{eq:Delta_matrix}
   \Delta \;=\; (\rho^{-1})^T \otimes \rho
   \quad \text{acting on}\quad \mathrm{vec}(X) \in \C^{\dim_W^2}.
\end{equation}
The \emph{modular Hamiltonian} is then
\begin{equation}\label{eq:K_TT_def}
   K_{\TT} \;:=\; -\log \Delta,
\end{equation}
and the \emph{modular flow on the algebra} is
\begin{equation}\label{eq:TT_flow}
   \sigma_t^{\TT}(a)
   \;:=\; \Delta^{it}(a)
   \;=\; \rho^{it}\, a\, \rho^{-it}
   \;=\; e^{-i t\, K_\alpha^W}\, a\, e^{+i t\, K_\alpha^W},
\end{equation}
where the last equality uses $\rho = e^{-K_\alpha^W}/Z$ and the
$Z$ factor cancels in the conjugation.

\subsection{The Tomita antilinear conjugation}\label{sec:J_TT_distinct}

A structurally important caution is that the Tomita modular
conjugation $J_{\TT}$ defined in~\eqref{eq:J_TT_def} is
\emph{categorically different} from the Connes KO-dim 3 charge
conjugation $J_{\mathrm{GV}}$ on the underlying Camporesi--Higuchi
spectral triple.

\begin{proposition}[$J_{\TT} \ne J_{\mathrm{GV}}$]\label{prop:J_TT_distinct}
The Tomita conjugation $J_{\TT}$ on $\Hilb_{\GNS}$ and the Connes
real structure $J_{\mathrm{GV}}$ on $\Hilb_{\nmax}$ are categorically
distinct antilinear operators with the following signature distinction:
\begin{equation}\label{eq:J_signatures}
   J_{\TT}^2 \;=\; +I, \qquad J_{\mathrm{GV}}^2 \;=\; -I.
\end{equation}
\end{proposition}

\begin{proof}
For $J_{\TT}^2 = +I$:\ applying~\eqref{eq:J_TT_def} twice gives
\begin{equation*}
   J_{\TT}^2(X)
   \;=\; J_{\TT}\bigl(\rho^{1/2}\, X^*\, \rho^{-1/2}\bigr)
   \;=\; \rho^{1/2}\,(\rho^{1/2}\, X^*\, \rho^{-1/2})^*\, \rho^{-1/2}
   \;=\; \rho^{1/2}\, \rho^{-1/2}\, X\, \rho^{1/2}\, \rho^{-1/2}
   \;=\; X,
\end{equation*}
where the second-to-last step uses $(\rho^{1/2})^* = \rho^{1/2}$ and
$(\rho^{-1/2})^* = \rho^{-1/2}$ (both are self-adjoint), and the last
step uses $\rho^{1/2} \rho^{-1/2} = I$ twice. For $J_{\mathrm{GV}}^2 = -I$:\
the Connes real structure on a KO-dim 3 spectral triple satisfies
$J^2 = -I$ by convention (cf.~Paper~32~\S~IV
\texttt{rem:axiom\_audit}).
\end{proof}

The two antilinear conjugations coexist on the same Hilbert spaces:\
$J_{\mathrm{GV}}$ on $\Hilb_{\nmax}$, $J_{\TT}$ on $\Hilb_{\GNS} \cong
B(\Hilb_W)$. Conflating them (e.g., by importing
\texttt{geovac.real\_structure.RealStructure} into the
modular-Hamiltonian construction) would silently produce algebraic
errors, since $J_{\TT}^2 - (-I) = 2I \ne 0$.

\begin{remark}[State-dependence vs.\ kinematic intrinsicness]
\label{rem:state_dependent}
$J_{\mathrm{GV}}$ is intrinsic to the spectral triple at construction
time:\ once $(\Acal, \Hilb, D, J, \gamma)$ is specified,
$J_{\mathrm{GV}}$ is fixed by the KO-dim convention. $J_{\TT}$ is
\emph{state-dependent}:\ each KMS state $\omega$ defines its own
$J_{\TT}^{(\omega)}$ via Tomita's theorem
on the GNS triple of $\omega$. The state-dependence is consistent
with the standard Tomita--Takesaki theorem~\cite{takesaki1970, tomita1967}.
\end{remark}

\subsection{Bit-exact period closure at the Tomita level}
\label{sec:tt_period}

\begin{theorem}[BW-$\gamma$ period closure at finite cutoff, {[SYMBOLIC]};
automatic, as Theorem~\ref{thm:bw_alpha}]
\label{thm:bw_gamma}
For every $\nmax \ge 1$ and every $a \in B(\Hilb_W)$, the
Tomita--Takesaki modular flow~\eqref{eq:TT_flow} satisfies the
period closure
\begin{equation}\label{eq:bw_gamma_closure}
   \sigma_{2\pi}^{\TT}(a) \;=\; a
\end{equation}
identically, with maximum periodicity residual
\begin{equation*}
   r_W^{\TT}(O)
   \;:=\; \Frob{\sigma_{2\pi}^{\TT}(O) - O}
\end{equation*}
scaling as $O(\sqrt{\dim_W} \cdot \varepsilon_{\mathrm{machine}})$ in
floating-point computation.
\end{theorem}

\begin{proof}
From~\eqref{eq:TT_flow}, $\sigma_{2\pi}^{\TT}(a) = e^{-i \cdot 2\pi
\cdot K_\alpha^W}\, a\, e^{+i \cdot 2\pi \cdot K_\alpha^W}$. By
Proposition~\ref{prop:integer_spectrum}, $K_\alpha^W$ has integer
spectrum, so $e^{\pm i \cdot 2\pi \cdot K_\alpha^W} = I$ identically.
Hence $\sigma_{2\pi}^{\TT}(a) = a$.
\end{proof}

\paragraph{Numerical verification.}
In floating-point computation, the modular operator $\Delta$ is built
on the GNS Hilbert--Schmidt space $\Hilb_{\GNS} \cong \C^{\dim_W^2}$
via the tensor-product form~\eqref{eq:Delta_matrix}, with each
$\rho^{\pm 1}$ assembled by matrix exponential of $\pm K_\alpha^W$ and
a numerical division by $Z$. The modular flow $\sigma_t^{\TT}$ is then
computed by acting $\Delta^{it}$ on each input operator $a$ written in
$\mathrm{vec}$ form, then reshaping. We have computed the maximum
periodicity residual $r_W^{\TT}(O)$ at $\nmax \in \{2, 3, 4, 5\}$;\
results are in Table~\ref{tab:bw_gamma_residuals}.

\begin{table}[ht]
\centering
\caption{Maximum BW-$\gamma$ Tomita-level periodicity residual
$\max_O r_W^{\TT}(O)$ at finite cutoff, computed via polar
decomposition on the GNS Hilbert--Schmidt space. The slightly larger
residual at $\nmax = 5$ (vs.~Table~\ref{tab:bw_alpha_residuals}) is
pure numerical conditioning from the
$\dim_W^2 \times \dim_W^2 = 4900 \times 4900$ matrix diagonalisation
of $\Delta$ with extreme eigenvalue ratios
$\lambda_i / \lambda_j$;\ no structural drift. $^{*}$Module label,
reproduced as computed;\ it records the automatic closure, not an
identification (the state is built from the posited generator). Source data:\
\texttt{debug/data/l1\_tighten\_tomita\_results.json}.}
\label{tab:bw_gamma_residuals}
\begin{tabular}{rrrcc}
\toprule
$\nmax$ & $\dim_W$ & $\dim \Hilb_{\GNS}$ &
$\max_O r_W^{\TT}(O)$ &
verdict\\
\midrule
2 & 8  & 64   & $1.09 \times 10^{-16}$ & UNIFIED\_STRONG$^{*}$ \\
3 & 20 & 400  & $1.07 \times 10^{-15}$ & UNIFIED\_STRONG$^{*}$ \\
4 & 40 & 1600 & $3.03 \times 10^{-15}$ & UNIFIED\_STRONG$^{*}$ \\
5 & 70 & 4900 & $4.79 \times 10^{-15}$ & UNIFIED\_STRONG$^{*}$ \\
\bottomrule
\end{tabular}
\end{table}

The $J_{\TT}^2 = +I$ identity of Proposition~\ref{prop:J_TT_distinct}
is verified numerically at every $\nmax$ to machine precision
(residuals $\le 7 \times 10^{-17}$):\ this distinguishes $J_{\TT}$
from $J_{\mathrm{GV}}$ at the genuine state-dependent level, not just
at the canonical $U = I$ sanity case.

\subsection{Spectrum of $K_{\TT}$}\label{sec:K_TT_spectrum}

\begin{proposition}[Spectrum of $K_{\TT}$ on $\Hilb_{\GNS}$]
\label{prop:K_TT_spectrum}
The eigenvalues of $K_{\TT} = -\log \Delta$ on $\Hilb_{\GNS}$ are
$\{n_j - n_k : 1 \le j, k \le \dim_W\}$, where $\{n_j\}$ is the integer
spectrum of $K_\alpha^W$. In particular:
\begin{enumerate}\setlength{\itemsep}{2pt}
\item[(a)] $\mathrm{Spec}(K_{\TT}) \subset \Z$ (integer spectrum on
$\Hilb_{\GNS}$);
\item[(b)] $K_{\TT}$ has at least $\dim_W$ zero eigenvalues (from the
diagonal $j = k$ pairs);
\item[(c)] The maximum eigenvalue of $K_{\TT}$ is $2 \max_j n_j =
2(2(\nmax - 1) + 1) = 2(2\nmax - 1)$;\ for $\nmax = 2$ this is $\le 4$,
and the maximum is exactly $2$ on the wedge (since $\max_W
\mathrm{two}\_m_j = +1$ at $\nmax = 2$).
\end{enumerate}
\end{proposition}

\begin{proof}
By~\eqref{eq:Delta_matrix}, $\Delta = (\rho^{-1})^T \otimes \rho$ in
the vec representation, with $\rho = e^{-K_\alpha^W}/Z$. The
eigenvalues of $\rho$ are $\{e^{-n_j}/Z\}_{j = 1}^{\dim_W}$, so the
eigenvalues of $\Delta$ are $\{e^{n_k} \cdot e^{-n_j}\}_{j, k} =
\{e^{n_k - n_j}\}_{j, k}$. Hence $K_{\TT} = -\log \Delta$ has
eigenvalues $\{n_j - n_k\}$, which are integer differences of the
integer spectrum of $K_\alpha^W$. The diagonal $j = k$ pairs contribute
$n_j - n_j = 0$, giving at least $\dim_W$ zero eigenvalues (more if
$K_\alpha^W$ has spectral degeneracies).
\end{proof}

The integer-spectrum property
of Proposition~\ref{prop:K_TT_spectrum}(a) is the structural reason
the bit-exact closure of Theorem~\ref{thm:bw_gamma} holds:\
$e^{i \cdot 2\pi \cdot K_{\TT}} = I$ on $\Hilb_{\GNS}$ by
$e^{i \cdot 2\pi \cdot m} = 1$ for integer $m$, so the modular flow
$\Delta^{it}$ acting at $t = 2\pi$ collapses to the identity, and
$\sigma_{2\pi}^{\TT}(a) = I \cdot a \cdot I = a$ for every $a$. This is
the same mechanism as Theorem~\ref{thm:bw_alpha} but realised at the
GNS-Hilbert--Schmidt level.

\subsection{State-space GH rate cross-check}\label{sec:propinquity_consistency}

The closure of Theorem~\ref{thm:bw_gamma} is a finite-cutoff identity,
not a propinquity limit statement. To confirm consistency with
Paper~38's qualitative-rate state-space-GH bound~\eqref{eq:paper38_main},
we compute the ratio
\begin{equation*}
   \frac{\max_O r_W^{\TT}(O)}{\gamma_{\nmax}^{(\mathrm{L2})}}
\end{equation*}
where $\gamma_{\nmax}^{(\mathrm{L2})}$ is the L2 mass-concentration
moment of the central spectral Fej\'er kernel on $\SU(2)$
(\cite[Lemma~L2]{paper38}):\
$\gamma_2^{(\mathrm{L2})} = 2.0746$, $\gamma_3^{(\mathrm{L2})} = 1.6101$,
$\gamma_4^{(\mathrm{L2})} = 1.3223$, $\gamma_5^{(\mathrm{L2})} = 1.1302$.
The ratios are $8.7 \times 10^{-17}$, $3.0 \times 10^{-16}$,
$7.1 \times 10^{-16}$, $1.4 \times 10^{-15}$:\ several orders of
magnitude below the qualitative-rate bound at every tested $\nmax$.
The comparison is a sanity check only and carries no content:\ the
closure is exact for the integer-spectrum reason, so its residual is
pure round-off and would be equally small for any integer-spectrum
generator, whatever the metric convergence rate.

% =====================================================================
\section{Unified closure}\label{sec:unified_closure}
% =====================================================================

The two constructions of Sections~\ref{sec:geometric_construction}
and~\ref{sec:tt_construction} are not the same operator on
$\Hilb_W$ versus $\Hilb_{\GNS}$ (the GNS construction lives on a
larger Hilbert space). The modular flows they generate on the
algebra $B(\Hilb_W)$ are the same flow up to time reversal, because
the state was built from $K_\alpha^W$;\ they are not structurally
independent realisations.  This section makes the relation precise.

\subsection{Flow conjugacy}\label{sec:flow_conjugacy}

\begin{theorem}[BW-$\alpha$ / BW-$\gamma$ flow conjugacy, {[SYMBOLIC]};
Tomita's theorem for a Gibbs state]
\label{thm:unified}
For every $a \in B(\Hilb_W)$ and every $t \in \R$,
\begin{equation}\label{eq:tt_alpha_conjugate}
   \sigma_t^{\TT}(a) \;=\; \sigma_{-t}^{\alpha}(a),
\end{equation}
where $\sigma_t^{\alpha}$ is the geometric BW-$\alpha$ flow
\eqref{eq:alpha_flow} and $\sigma_t^{\TT}$ is the Tomita--Takesaki
BW-$\gamma$ flow~\eqref{eq:TT_flow}.
\end{theorem}

\begin{proof}
By~\eqref{eq:TT_flow}, $\sigma_t^{\TT}(a) = e^{-it K_\alpha^W}\, a\,
e^{+it K_\alpha^W}$. By~\eqref{eq:alpha_flow}, $\sigma_{-t}^{\alpha}(a) =
e^{-it K_\alpha^W}\, a\, e^{+it K_\alpha^W}$. The two expressions are
literally identical.
\end{proof}

\paragraph{Consequence at the period.}
At $t = 2\pi$, equation~\eqref{eq:tt_alpha_conjugate} reads
$\sigma_{2\pi}^{\TT}(a) = \sigma_{-2\pi}^{\alpha}(a) =
\sigma_{2\pi}^{\alpha}(a)^{-1}$, where the last equality uses that the
flow $\sigma_t^{\alpha}$ is a one-parameter group with the inverse at
$t = 2\pi$ being the flow at $t = -2\pi$. But by
Theorem~\ref{thm:bw_alpha}, $\sigma_{2\pi}^{\alpha}(a) = a = a^{-1}$
trivially (this is an operator equation, not a group inverse), so
$\sigma_{2\pi}^{\TT}(a) = a$ as well, consistent with
Theorem~\ref{thm:bw_gamma}. The two closures are the \emph{same}
closure under the conjugacy \eqref{eq:tt_alpha_conjugate}, which is
why their agreement carries no independent evidential weight.

\paragraph{Numerical confirmation at general $t$.}
At a general non-period $t$ (e.g., $t = 1$), the two flows
$\sigma_t^{\TT}$ and $\sigma_t^{\alpha}$ are distinct (the conjugacy
\eqref{eq:tt_alpha_conjugate} relates $\sigma_t^{\TT}$ to
$\sigma_{-t}^{\alpha}$, not to $\sigma_t^{\alpha}$). We have verified
this numerically:\ the average cross-flow difference
$\langle \Frob{\sigma_1^{\TT}(O) - \sigma_1^{\alpha}(O)} \rangle_O$
over the same five test operators is $\approx 0.16$ at $\nmax = 2$,
clearly non-zero. The sign-flipped comparison
$\langle \Frob{\sigma_1^{\TT}(O) - \sigma_{-1}^{\alpha}(O)} \rangle_O$
gives $\approx 4 \times 10^{-17}$, machine precision. The conjugacy
\eqref{eq:tt_alpha_conjugate} is therefore verified bit-exactly at
general $t$, not just at the period.

\subsection{The posited local Hamiltonian}
\label{sec:derived_hamiltonian}

The unified closure rests on the deliberate choice
$H_{\mathrm{local}} = K_\alpha^W / \beta$
(equation~\eqref{eq:H_local_def}) for the local Hamiltonian
generating the wedge KMS state. This is a construction-level choice,
not a derivation:\ we wrote down the wedge KMS state as $\rho_W =
e^{-K_\alpha^W}/Z$ rather than (say) $\rho_W = e^{-\beta D_W}/Z$ for
the wedge-restricted Camporesi--Higuchi Dirac. The motivation is
structural (we want the modular flow to close with period $2\pi$, as
the Bisognano--Wichmann flow does), but it is \emph{a} choice of local
Hamiltonian, not the unique choice.  The generator $K_\alpha^W = 2J_z$
is \emph{posited}, and $\rho_W = e^{-K_\alpha^W}/Z$ is its unit-$\beta$
Gibbs state, not the Bisognano--Wichmann-normalised state, which (KMS
parameter $2\pi$ in the rotation angle) would be
$e^{-2\pi J_z}/Z = e^{-\pi K_\alpha^W}/Z$ (Paper~50, remark on the
normalisation of $K_\alpha$ and $\rho_W$).

Two remarks place the choice;\ neither derives it:

\paragraph{(I) Formal resemblance to the Wightman BW vacuum.}
In the continuum Bisognano--Wichmann theorem, the local Hamiltonian
on the Rindler wedge is the boost generator $K_{\mathrm{boost}}$, and
the Wightman vacuum restricted to the wedge algebra is the Gibbs state
$\rho_W = e^{-2\pi K_{\mathrm{boost}}}/Z$ at inverse temperature
$\beta = 2\pi$. Our choice $H_{\mathrm{local}} = K_\alpha^W / \beta$
copies this \emph{form} (not its normalisation, see above):\
$\beta H_{\mathrm{local}} = K_\alpha^W$, with
the spatial rotation generator around the Hopf-base axis standing in
for the boost, and $\rho_W = e^{-K_\alpha^W}/Z$.  In the continuum the
Gibbs form is a \emph{theorem} about the vacuum;\ here it is the
definition of the state, so the resemblance is by design and is not
evidence that $\rho_W$ is the restriction of any distinguished global
state.

\paragraph{(II) The intrinsic Dirac lifts the period to $-I$.}
If we had chosen $H_{\mathrm{local}} = D_W$ (the wedge-restricted
Camporesi--Higuchi Dirac), whose eigenvalues are $\chi \cdot (n + 1/2)$
(half-integer on the spinor basis), then the \emph{operator-level}
period lift $e^{i \cdot 2\pi \cdot D_W} = -I$ --- the global $-1$ of the
spinor double cover --- whereas the posited $K_\alpha^W$ has
odd-integer spectrum and $e^{i \cdot 2\pi \cdot K_\alpha^W} = +I$, the
identity operator. We emphasise the precise sense of the distinction:\
the modular flow $\sigma_t(\cdot) = U_t (\cdot) U_t^{-1}$ acts by
\emph{conjugation}, in which a global scalar phase cancels, so the
flow-level period closure $\sigma_{2\pi}(O) = O$ still holds with $D_W$
(the $-I$ cancels:\ $(-I)\, O\, (-I) = O$;
we verified this bit-exactly, residual $<10^{-13}$ at $n_{\max} =
2,3,4$;\ \texttt{tests/test\_modular\_hamiltonian.py}). The distinction
is therefore \emph{not} at the flow level but
at the operator level:\ only $K_\alpha^W$ lifts the $\beta = 2\pi$
period to the identity operator $+I$; $D_W$ lifts it to the
double-cover $-I$.

This is an elementary observation (integer versus half-integer
spectrum), not a selection theorem. The framework's intrinsic
Dirac $\DCH$ is the canonical spectral-action object on the
Camporesi--Higuchi triple, and one might expect it to play the
role of the local Hamiltonian for a natural KMS state on the
truncated operator system. The above shows only that, while either
choice generates a $2\pi$-periodic conjugation flow, the posited
generator $K_\alpha^W$ lifts the period to $+I$ and the
wedge-restricted Dirac $D_W$ to the spinor double-cover sign $-I$;\
which (if either) is the modular generator of a Camporesi--Higuchi
wedge state is the [OPEN] question O0. The two play
complementary roles:\
$\DCH$ is the spectral-action object (the Connes--Chamseddine input
for the heat-kernel expansion); $K_\alpha^W$ is the
wedge-modular-Hamiltonian object (the Tomita--Takesaki input for the
modular-flow construction). On the round $\sthree$ they happen to
agree up to an overall scale (both are diagonal in the spinor basis
with simple spectral structure), but only one of them has the integer
spectrum that lifts the $\beta = 2\pi$ period to the identity operator
(the same integer-spectrum / $e^{2\pi i K} = +I$ structure
that underlies the compact-boost reading of
Section~\ref{sec:conclusion}).

\begin{remark}[A speculative reading]\label{rem:speculation}
We note, but do not pursue here, the possibility that the
distinction between the spectral-action Dirac $\DCH$ and the
modular-Hamiltonian generator $K_\alpha^W$ is structurally significant
for the framework's interpretation. In the spectral-action paradigm
(Chamseddine--Connes~\cite{chamseddine_connes2010}), the Dirac is the
universal object from which all physics is read; in the
thermal-time paradigm (Connes--Rovelli~\cite{connes_rovelli1994}), the
modular Hamiltonian of an algebraic state is the physical-time generator.
On the round $\sthree$ Camporesi--Higuchi triple at the wedge KMS
state $\rho_W = e^{-K_\alpha^W}/Z$, these are not the same operator.
The structural status of this distinction --- whether it is a
peculiarity of the round $\sthree$ truncation or a deeper feature of
the framework's spectral content --- is open. We flag it as one of
the open questions in Section~\ref{sec:conclusion}.
\end{remark}

\subsection{Cross-witness collapse at the algebra-action level}
\label{sec:cross_witness_algebra}

The witness agreement mentioned in
Section~\ref{sec:four_witness_unit} holds by construction at the
algebra-action level. With $H_{\mathrm{local}} = K_\alpha^W / \beta$,
the density matrix is
\begin{equation}\label{eq:rho_W_beta_indep}
   \rho_W \;=\; \frac{e^{-\beta H_{\mathrm{local}}}}{Z}
            \;=\; \frac{e^{-K_\alpha^W}}{Z}
\end{equation}
\emph{independent of $\beta$}. The Tomita modular operator $\Delta$
constructed from $\rho_W$ via~\eqref{eq:Delta_matrix} and the modular
Hamiltonian $K_{\TT} = -\log \Delta$ are therefore also
$\beta$-independent. All four witness instantiations (BW, HH, Sew,
Unruh) therefore produce the \emph{same} modular operator $\Delta$ and
the same $K_{\TT}$:\ they are the same $K_\alpha^W$ with $\kappa_g$
relabelled, and $\kappa_g$ drops out by construction.  This is not a
finite-cutoff statement about the four witnesses.

We have verified this numerically at $\nmax \in \{2, 3, 4, 5\}$ across
the six witness instantiations (BW witness at $\kappa_g = 1$;
Hartle--Hawking at $M = 1$, $\kappa_g = 1/4$; Hartle--Hawking at
$M = 2$, $\kappa_g = 1/8$; Sewell at $M = 1$, $\kappa_g = 1/4$;
Unruh at $a = 1$, $\kappa_g = 1$; Unruh at $a = 2$, $\kappa_g = 2$).
The cross-witness consistency residual is exactly $0$ at every tested
$\nmax$ between any pair of witnesses, as it must be:\ the six runs
compute the same matrix, so the check confirms the bookkeeping, not a
physical coincidence.

% =====================================================================
\section{Four witnesses, one construction (by construction)}
\label{sec:four_witness_collapse}
% =====================================================================

The numerical content of the witness agreement is summarised in
Table~\ref{tab:cross_witness} and its (elementary) structural content
is the following corollary.

\begin{table}[ht]
\centering
\caption{Cross-witness consistency at the BW-$\alpha$ level
(left columns) and at the BW-$\gamma$ Tomita level (right columns).
Cross-consistency residual is the difference between the period
residual of each witness and the BW-witness period residual at the
same $\nmax$:\ exactly $0$ across all six witnesses
at every tested $\nmax$. Source data:\
\texttt{debug/data/l1\_modular\_hamiltonian\_results.json},
\texttt{debug/data/l1\_tighten\_tomita\_results.json}.}
\label{tab:cross_witness}
\renewcommand{\arraystretch}{1.1}
\setlength{\tabcolsep}{4pt}
\begin{tabular}{rcccccc}
\toprule
$\nmax$ & $\BW$ & $\HH_{M=1}$ & $\HH_{M=2}$ & $\Sew_{M=1}$ & $\Unruh_{a=1}$ & $\Unruh_{a=2}$ \\
& ($\kappa_g\!=\!1$) & ($\kappa_g\!=\!1/4$) & ($\kappa_g\!=\!1/8$) & ($\kappa_g\!=\!1/4$) & ($\kappa_g\!=\!1$) & ($\kappa_g\!=\!2$) \\
\midrule
\multicolumn{7}{l}{\emph{BW-$\alpha$ residual (Table~\ref{tab:bw_alpha_residuals}):}} \\
2 & 1.80e-16 & 1.80e-16 & 1.80e-16 & 1.80e-16 & 1.80e-16 & 1.80e-16 \\
3 & 4.76e-16 & 4.76e-16 & 4.76e-16 & 4.76e-16 & 4.76e-16 & 4.76e-16 \\
4 & 9.33e-16 & 9.33e-16 & 9.33e-16 & 9.33e-16 & 9.33e-16 & 9.33e-16 \\
5 & 1.59e-15 & 1.59e-15 & 1.59e-15 & 1.59e-15 & 1.59e-15 & 1.59e-15 \\
\midrule
\multicolumn{7}{l}{\emph{BW-$\gamma$ Tomita residual (Table~\ref{tab:bw_gamma_residuals}):}} \\
2 & 1.09e-16 & 1.09e-16 & 1.09e-16 & 1.09e-16 & 1.09e-16 & 1.09e-16 \\
3 & 1.07e-15 & 1.07e-15 & 1.07e-15 & 1.07e-15 & 1.07e-15 & 1.07e-15 \\
4 & 3.03e-15 & 3.03e-15 & 3.03e-15 & 3.03e-15 & 3.03e-15 & 3.03e-15 \\
5 & 4.79e-15 & 4.79e-15 & 4.79e-15 & 4.79e-15 & 4.79e-15 & 4.79e-15 \\
\bottomrule
\end{tabular}
\end{table}

\begin{corollary}[Single operator-system construction for the four
witnesses, {[SYMBOLIC]}; by construction]\label{cor:single_construction}
Under the unit normalisation U-1 and the wedge KMS state choice
$\rho_W = e^{-K_\alpha^W}/Z$, the operator-system-level realisations of
the Hartle--Hawking, Sewell, Bisognano--Wichmann, and Unruh witnesses
on the truncated Camporesi--Higuchi triple $\Tcal_{\nmax}$ collapse to
a single construction. The modular operator $\Delta$ and the modular
Hamiltonian $K_{\TT}$ are bit-identical across the six witness
instantiations. The witness-specific physics enters only through the
parameterisation of $\kappa_g$, which is invisible at the
algebra-action level and at the period closure $\sigma_{2\pi}(O) = O$.
\end{corollary}

\begin{proof}
The density matrix $\rho_W = e^{-K_\alpha^W}/Z$ is
$\beta$-independent (Section~\ref{sec:cross_witness_algebra}). The
Tomita modular operator $\Delta$ on the GNS Hilbert--Schmidt space is
constructed entirely from $\rho_W$ via~\eqref{eq:Delta_matrix},
and the modular Hamiltonian $K_{\TT} = -\log \Delta$ depends only on
$\Delta$. Both are therefore $\beta$-independent.
$\beta$-independence at the algebra-action level immediately implies
identical period closure at $t = 2\pi$ for all six witness
instantiations.
\end{proof}

\paragraph{Scope of the corollary.}
The corollary is true because the six instantiations are the same
construction with a relabelled parameter;\ it is not a finite-cutoff
theorem about Hartle--Hawking, Sewell, Bisognano--Wichmann, or Unruh.
The witness-specific physical content (mass of the black hole, proper
acceleration of the Rindler observer) enters only through the label
$\kappa_g$, which the construction never uses.  The $2\pi$ of the
period closure comes from the integer spectrum of the posited $2J_z$;\
it is \emph{not} shown here to be the same $2\pi$ as the M1 volume
factor of the Paper~38 rate constant~\eqref{eq:paper38_main}, and an
earlier version's reading of the closure as the operator-system
realisation of the M1 sub-mechanism is withdrawn (2026-09-28).
The metric-functional-level reading
(\cite[Rem.~rem:bisognano\_wichmann\_reading]{paper32}, Sprint TD
Track 4, Sprint Unruh-pendant), which matches the framework's $2\pi$
on a temporal $S^1_\tau$ with the modular-flow / boost-orbit period
of the four witnesses via the published Wick-rotation chain, is
unaffected and is \emph{not} lifted by the present paper:\ the
dependence on the Wick-rotation prescription remains, pending the
[OPEN] derivation of $K$ from a wedge-restricted Camporesi--Higuchi
state.

\begin{remark}[Hilbert--Schmidt orthogonality of $H_{\mathrm{local}}$ and
$D_W$]\label{rem:pythagorean_underlies_collapse}
For every witness $w$ the posited local Hamiltonian
$H_{\mathrm{local}}^{(w)} = K_\alpha^W/\beta^{(w)}$ and the wedge-restricted
Camporesi--Higuchi Dirac $D_W^{\mathrm{GV}}$ are Hilbert--Schmidt
orthogonal, $\Tr\big((H_{\mathrm{local}}^{(w)})^\dagger D_W^{\mathrm{GV}}\big) = 0$,
so that $\Frob{H_{\mathrm{local}}^{(w)} - D_W^{\mathrm{GV}}}^2 =
(\kappa_g^{(w)})^2 S(\nmax)/(4\pi^2) + D(\nmax)$ with $S, D$ polynomials in
$\nmax$.  This is a chirality-parity identity, proved in
Theorem~\ref{thm:pythagorean_formal} and
Corollary~\ref{cor:pythagorean_orthogonality}
(\S\ref{sec:pythagorean_orthogonality}):\ $H_{\mathrm{local}}^{(w)}$ is
diagonal and independent of the chirality $\chi$, while $D_W^{\mathrm{GV}}$
is chirality-diagonal with $\chi$-paired eigenvalues of opposite sign, so
the pairs cancel in the trace.  (An earlier version described
$D_W^{\mathrm{GV}}$ as off-diagonal in the wedge basis;\ it is diagonal and
sign-paired.)  The identity is $\beta$-independent, the $1/\pi^2$ of the
residual being the square of the posited normalisation
$\beta = 2\pi/\kappa_g$.  It does not ``underlie'' the six-witness agreement
of Corollary~\ref{cor:single_construction}, which holds by construction,
and does not bear on the period closure, which is automatic for the
integer-spectrum reason of Remark~\ref{rem:bw_alpha_automatic}.
\end{remark}

% =====================================================================
\section{Related work and scope}\label{sec:related_work}
% =====================================================================

\paragraph{Connes--Rovelli thermal time hypothesis.}
The closest published analog of the present construction is the
Connes--Rovelli thermal time hypothesis~\cite{connes_rovelli1994},
which identifies the Tomita modular flow of a general-covariant von
Neumann algebra with the physical-time evolution. The present paper
writes the Tomita machinery out on a finite-dimensional Type
I$_{\dim_W}$ factor obtained as the wedge-restricted
Connes--van~Suijlekom~\cite{connes_vs2021} truncation of the round
$\sthree$ Camporesi--Higuchi spectral triple, for a state that is
itself defined from the chosen generator.  The Connes--Rovelli
construction starts from a state and reads off its flow;\ the present
construction starts from the flow and builds the state, so it
illustrates the machinery rather than testing the hypothesis.

\paragraph{Bisognano--Wichmann tradition and continuum modular
Hamiltonians.}
The original Bisognano--Wichmann
papers~\cite{bisognano_wichmann1975, bisognano_wichmann1976} establish
the modular-flow-as-boost identification for the Rindler wedge of
Minkowski spacetime, with the analytic period $\sigma_{i \cdot 2\pi} =
\mathrm{id}$ as a corollary of the Tomita--Takesaki KMS
condition~\cite{takesaki1970, tomita1967} at $\beta = 2\pi$ in
rapidity units. The continuum entanglement-Hamiltonian construction
for a half-space in free CFT was developed by Casini and
Huerta~\cite{casini_huerta2009}, and extended to spherical regions of
$S^d$ in CFT by Casini--Huerta--Myers~\cite{casini_huerta_myers2011}.
The present construction is a formal analog of the
Casini--Huerta--Myers spherical-region modular Hamiltonian on the
truncated Camporesi--Higuchi triple over the round $\sthree$, with
the integer spectrum of the posited $K_\alpha^W$ standing in for the
canonical $2\pi$-period of the continuum modular flow;\ unlike
Casini--Huerta--Myers, it does not derive the modular Hamiltonian from
a state.

\paragraph{Lattice realisations.}
Zhang~\emph{et~al.}~\cite{zhu_casini2020} verified the
Bisognano--Wichmann form of the modular Hamiltonian on critical
quantum spin chains, finding the continuum BW prediction to hold to
good approximation on finite lattices. The present paper differs in
three respects:\ (i)~the underlying object is a continuum spectral
triple on a non-flat compact $\sthree$ background, not a lattice
discretisation of physical space; (ii)~the finite-dimensional
structure comes from the spectral truncation of $\DCH$ via the
Connes--van~Suijlekom framework, not from a real-space discretisation;
(iii)~most importantly, Zhang et al.\ \emph{extract} the modular
Hamiltonian from a ground state and compare it with the BW form,
which is a genuine test, whereas the present paper posits the
generator, so its period closure is automatic.  The analog of their
test on the truncated triple is the [OPEN] question O0.

\paragraph{Algebraic QFT on curved backgrounds.}
The standard algebraic-QFT references for modular structures on
Riemannian backgrounds are
Strohmaier~\cite{strohmaier2006} (Reeh--Schlieder on stationary
spacetimes) and Verch~\cite{verch2001} (generally covariant
spin-statistics and modular structure on curved
backgrounds). These address the modular-theoretic content
of axiomatic QFT on curved backgrounds. The present construction is
spectral-triple-internal:\ it does not invoke a Wightman field algebra
or a nuclearity argument, but rather works directly with the
$\Op_{\nmax} \subset B(\Hilb_{\nmax})$ truncation of the
$C^\infty(\sthree) \subset B(L^2(\sthree, \Sigma))$ continuum algebra
and the associated Hilbert--Schmidt GNS construction.

\paragraph{Lorentzian signature and Krein-space spectral triples.}
The proper Lorentzian extension of the present construction would
require a Krein-space spectral triple at Lorentzian signature (one time
and three space dimensions), for which
the canonical reference is
Bizi--Brouder--Besnard~2018~\cite{bizi_brouder_besnard2018}. The
Bizi--Brouder--Besnard classification organises spectral triples of
metric signature $(t, s)$ ($t$ time and $s$ space dimensions) by pairs
$(m, n) = (t + s,\, t - s) \bmod 8$ (their p.~3;\ all BBB table,
equation and page numbers in this paper follow arXiv:1611.07062 v2),
where $m$ alone is the metric dimension, and the underlying Connes axiom
signs determine the Krein-space structure.
The Riemannian Camporesi--Higuchi triple has $(t, s) = (0, 3)$ (three
spatial directions), which gives $(m, n) = (3, -3 \bmod 8) = (3, 5)$
--- this is BBB's formula extended formally to odd $m$ (their tables
cover even $m, n$ only);\ Connes' KO-dimension~3
corresponds to it;\ the
Lorentzian extension, one time and three space dimensions, has
$(t, s) = (1, 3)$, which gives $(m, n) = (4, -2 \bmod 8) = (4, 6)$, their
Table~3 entry $(4, 6)$.  (Elsewhere in this paper a bare signature pair
$(3, 0)$ or $(3, 1)$ is written in the order $(s, t)$, space first;\ it is
never a BBB $(m, n)$ label.) The Bizi--Brouder--Besnard signs there (their
Table~1 and eq.~(3)) make $J$ anticommute with the chirality, which leads to the
present author's Sprint L0 prediction that the M3 sub-mechanism of
the master Mellin engine trivialises at $(t, s) = (1, 3)$ on chirality-symmetric
spectrum (\texttt{debug/lorentzian\_l0\_audit\_memo.md} \S 3).
The Krein-space lift and the replication of the present modular
construction at $(t, s) = (1, 3)$ are carried out at finite cutoff in
\S\ref{sec:lorentzian_closure};\ the M3 prediction there turns out to be
convention-dependent (\S\ref{sec:m3_convention}).  A continuum
Lorentzian analog of the Latr\'emoli\`ere propinquity remains absent from
the literature.

\paragraph{Twisted-spectral-triple emergence.}
The recent work of
Nieuviarts~\cite{nieuviarts2025a, nieuviarts2024, nieuviarts2025b_proceedings}
proposes a twisted-spectral-triple morphism that constructs a
pseudo-Riemannian spectral triple from a Riemannian almost-commutative
twisted spectral triple, presented as an algebraic alternative to
Wick rotation. The construction does not apply to the present
$\sthree$ case because of the structural odd-dimensional restriction
(\S~\ref{sec:nieuviarts_footnote}). The Lorentzian-extension route for
the present paper therefore runs through the direct
Bizi--Brouder--Besnard~2018 Krein-space lift, not through a twist
morphism.

\paragraph{Other NCG-side gaps the present paper does not close.}
A genuine Lorentzian propinquity (in the sense of an analog of the
Latr\'emoli\`ere quantum Gromov--Hausdorff propinquity for
Krein-space metric spectral triples) is open as of May 2026.
Toyota~\cite{toyota2023}, Farsi--Latr\'emoli\`ere~\cite{farsi_latremoliere2024,
farsi_latremoliere2025}, Rieffel~\cite{ucp_maps_2024},
and Hekkelman--McDonald~\cite{hekkelman_mcdonald2024b} all extend
the metric-spectral-triple propinquity machinery in various directions
(discrete groups of polynomial growth, collapse with abelian factors,
analytic paths of Riemannian metrics, noncommutative integral on
truncations, Berezin--Toeplitz on $S^2$), but none addresses
Krein-space signature change. Constructing the Lorentzian propinquity
is original NCG-mathematics work, not algorithmic.

% =====================================================================
\section{Conclusion and open questions}\label{sec:conclusion}
% =====================================================================

We have set up, on the truncated Camporesi--Higuchi spectral triple
$\Tcal_{\nmax}$, a hemispheric wedge $P_W$ aligned with the Hopf-base
axis, a wedge Gibbs state built from the posited generator
$K_\alpha^W = \mathrm{diag}(2m_j) = 2J_z$, and the corresponding
GNS/Tomita construction.  What holds is elementary and [SYMBOLIC]:\
the period closure $\sigma_{2\pi}(O) = O$
(Theorems~\ref{thm:bw_alpha}, \ref{thm:bw_gamma}), automatic for any
integer-spectrum generator and any $O$;\ the conjugacy
$\sigma_t^{\TT}(O) = \sigma_{-t}^{\alpha}(O)$
(Theorem~\ref{thm:unified}), Tomita's theorem for a Gibbs state built
from $K_\alpha^W$;\ and the agreement of the six witness
instantiations (Corollary~\ref{cor:single_construction}), which share
one generator and differ only in the label $\kappa_g$.  All are
reproduced bit-exactly in floating point at $\nmax \in \{2, 3, 4, 5\}$.
We also record the operator-level contrast between $2J_z$
($e^{2\pi i K} = +I$) and the wedge Dirac ($e^{2\pi i D_W} = -I$).

The same holds on the signature-$(3, 1)$ Krein carrier of
\S\ref{sec:lorentzian_closure}, where the period closure is again
automatic and signature-blind;\ that section also records the Krein axioms
and Krein-self-adjointness of the Lorentzian Dirac, the four
Bizi--Brouder--Besnard-predicted sign relations at $(4, 6)$, the failure
of their universal chirality anticommutation on the truthful
Camporesi--Higuchi Dirac, and the chirality-parity Hilbert--Schmidt
orthogonality of $H_{\mathrm{local}}$ and the wedge Dirac, whose residual
carries $1/\pi^2$ only through the posited normalisation.

\emph{What is withdrawn} (2026-09-28).\ \ An earlier version of this
paper read these closures as lifting the four-witness Wick-rotation
theorem from a structural correspondence to an identification at the
operator-system level at every finite cutoff.  That reading is
withdrawn:\ the generator is posited rather than derived, the Tomita
step is circular, and the witness agreement holds by construction, so
nothing here identifies a Bisognano--Wichmann modular flow at finite
cutoff.  \emph{What stands:}\ the continuum four-witness statement
($\beta = 2\pi/\kappa_g$) as prior art, and the elementary facts above.

\paragraph{Open questions.}

\textbf{(O0) Derive the modular Hamiltonian from a Camporesi--Higuchi
state} [OPEN].  The falsifier the corpus named for this construction
--- compute $K$ on $\Tcal_{\nmax}$ as the modular Hamiltonian of a
wedge-restricted Camporesi--Higuchi state (for instance the
restriction to $P_W \Op_{\nmax} P_W$ of a global state defined from
$\DCH$ alone), and only then test $\sigma_{2\pi} = \mathrm{id}$ and
compare $K$ with $2J_z$ --- has not been run.  It is the finite-cutoff
counterpart of the lattice test of Zhang~\emph{et~al.}~\cite{zhu_casini2020},
and it is the question on which any finite-cutoff four-witness claim
would rest.

\textbf{(O1) Lorentzian extension at signature $(t, s) = (1, 3)$}
[finite-cutoff part carried out;\ continuum part descoped].
As drafted, O1 proposed to replicate the present operator-system
construction on a Bizi--Brouder--Besnard Krein-space spectral triple
at BBB $(m, n) = (4, 6)$, to verify the Sprint L0 prediction that M3
trivialises but M1 stays under the BBB signs (their Table~1 and
eq.~(3)), and to produce a Lorentzian KMS--Tomita--Takesaki theorem
inside the framework.  The replication was carried out at finite cutoff
(\S\ref{sec:lorentzian_closure});\ the M3 prediction was falsified
under the $n_{\mathrm{fock}}$-parity convention
(\S\ref{sec:m3_convention});\ the Krein period closures are automatic
for the posited generator and yield no Lorentzian KMS--Tomita--Takesaki
theorem;\ and the continuum-propinquity leg was descoped
(Paper~45~\cite{paper45}).  The progress record below is kept as
history.

\emph{Progress update (Sprints L2-B/C/D, 2026-05-16).}
The Hilbert-space and Lorentzian-Dirac levels of the $(3, 1)$
extension have landed:\ Sprint L2-B
(\texttt{geovac/krein\_space\_construction.py}) constructs the
Krein space $\mathcal{K}_{\nmax, N_t} = \mathcal{H}_{\rm GV}^{\nmax}
\otimes L^2(\mathbb{R}_t)_{\rm cutoff}$ with fundamental symmetry
$J = \gamma^0$ in the Peskin--Schroeder chiral basis, verifying all
Krein axioms ($J^2 = +I$, $J^* J = I$, $J^* = J$,
$\mathcal{K} = \mathcal{K}^+ \oplus \mathcal{K}^-$) bit-exactly at
$(\nmax, N_t) \in \{1, 2, 3\} \times \{1, 11, 21\}$ and passing the
load-bearing Riemannian-limit recovery at $N_t = 1$ bit-identically;
Sprint L2-C (\texttt{geovac/lorentzian\_dirac.py}) constructs
$D_L = i \cdot [\gamma^0 \otimes \partial_t + D_{\rm GV} \otimes I_{N_t}]$
via van~den~Dungen Proposition~4.1 with $\partial_t$ centered FD plus
Dirichlet zero BC (structurally forced for anti-Hermiticity and
Krein-self-adjointness), verifying Krein-self-adjointness
$D_L^\times = D_L$ and Riemannian-limit recovery $D_L|_{N_t = 1} =
i D_{\rm GV}$ bit-exactly across the same panel.  Sprint L2-D
(\texttt{geovac/connes\_axiom\_audit\_31.py}) verifies the four
BBB-predicted sign relations at $(m, n) = (4, 6)$ bit-exactly
($J_L^2 = +I$, $\{J_L, \gamma^5\} = 0$, $\{J_L, \eta\} = 0$,
$J_L D_L = +D_L J_L$;\ Theorem~\ref{thm:connes_audit}), while the
order-zero and order-one conditions fail at finite cutoff
(Remark~\ref{rem:order_zero_one}) and the universal axiom
$\{\gamma^5, D_L\} = 0$ is a
\emph{structural finding}, failing on the truthful Camporesi--Higuchi
Dirac (which is chirality-diagonal in the chiral basis and therefore
commutes with $\gamma^5$;\ \S\ref{sec:bbb_axiom_finding}).  This is structurally
analogous to the Riemannian-side R3.5/present-paper trade between
truthful CH (correct for axioms requiring $JD = +DJ$) and offdiag CH
(correct for SDP-bounding Connes-distance constructions).  The
construction does \emph{not} re-open WH1's commutative-factor convergence result
(Paper~38~\cite{paper38}); the Krein construction is an extension on
top of the Riemannian foundation, not a re-test of it.  An
operator-system-level identification of the four-witness theorem at
Lorentzian signature (which would first require the Riemannian [OPEN]
question O0) was the named Sprint L2-E target, still open beyond
sprint scale due to the absence of a published Lorentzian
propinquity.  O1 status remains OPEN \emph{at the continuum /
non-compact level}; the Hilbert-space and Dirac-operator levels are now
closed at finite cutoff.  On the \emph{metric}-level finite-cutoff
question (Paper~45~\S\,open~Q1) only the compactness half is established
[SYMBOLIC]:\ the posited generator $K_\alpha = 2J_z$ is compact (integer
spectrum, $e^{2\pi i K} = I$).  The further reading has two legs,
priced separately.  (a)~That the finite-cutoff quantum metric is
Euclidean, with the genuine reverse-triangle signature confined to the
non-compact continuum limit, is [ARGUED] and \emph{generator-independent}:\
it does not depend on which generator a wedge state selects.
(b)~That the Lorentzian signature is the Wick rotation of the KMS circle,
a convention, is [ARGUED] and \emph{conditional on O0}:\ it applies to
the Bisognano--Wichmann generator only if $K$ is derived from a
wedge-restricted Camporesi--Higuchi state, which is open.

\textbf{(O2) Non-tracial wedge states.}
The closure rests on the Gibbs structure
$\rho_W = e^{-K_\alpha^W}/Z$ on a finite-dimensional Type I$_{\dim_W}$
factor. For non-tracial wedge states (e.g., excited states above the
BW vacuum, mixed states with non-Gibbs structure), the polar
decomposition $S = J_{\TT} \Delta^{1/2}$ still goes through, but the
integer-spectrum property of $K_{\TT}$ is not guaranteed, and the
bit-exact closure may degrade to a soft propinquity-rate-controlled
closure of the type Paper~38 supplies. The structural status of the
non-tracial case is open.

\textbf{(O3) Spectral-action Dirac vs.\ modular-Hamiltonian generator.}
Section~\ref{sec:derived_hamiltonian} flagged that the framework's
intrinsic Camporesi--Higuchi Dirac $\DCH$ lifts the $\beta = 2\pi$
period to $-I$, whereas the posited $K_\alpha^W$ lifts it to $+I$.
Which generator (if either) a Camporesi--Higuchi wedge state actually
selects is O0;\ the present paper chooses $K_\alpha^W$ rather than
deriving it. In
general the spectral-action $\DCH$ and the modular-Hamiltonian
$K_\alpha^W$ may not even commute. The structural status of this
distinction in extensions to other compact-Lie-group spectral triples
(per Paper~40~\cite{paper40_unified}) is open and an interesting
direction for follow-up.

\emph{Refinement (Sprint L2-E, 2026-05-16; \S~\ref{sec:lorentzian_closure} below).}\ \ %
The O3 question of signature-dependence is now \emph{partially}
closed:\ the $H_{\mathrm{local}} \ne D_W$ distinction is
\textbf{signature-INDEPENDENT at the Riemannian limit} and refined
upward at $N_t > 1$ by temporal-derivative content.  At $N_t = 1$
on the Lorentzian Krein space, $\|H_{\mathrm{local}} - D_L^W\|_F$
equals the Riemannian-side residual $\|H_{\mathrm{local}} - D_{\mathrm{GV}}^W\|_F$
\emph{bit-exact} at every tested $n_{\max} \in \{1, 2, 3\}$ (because
$D_L = iD_{\mathrm{GV}}$ at $N_t = 1$ and $H_{\mathrm{local}}$ is
real).  This is a property of the posited construction, not evidence
about the modular content of a Camporesi--Higuchi state.  At $N_t > 1$
the Lorentzian Dirac
$D_L = i(\gamma^0 \otimes \partial_t + D_{\mathrm{GV}} \otimes I)$
has temporal-derivative content that $K_\alpha^W / \beta$ cannot
capture; the gap between the two operators widens with temporal
cutoff.  Which generator a Camporesi--Higuchi wedge state selects
remains O0.

\emph{Further refinement (Hilbert--Schmidt orthogonality;\
\S\ref{sec:pythagorean_orthogonality}).}\ \ %
The signature-independent residual is Pythagorean:\
$\langle H_{\mathrm{local}}, D_W^L \rangle_{\mathrm{HS}} = 0$ as a
chirality-parity identity (Theorem~\ref{thm:pythagorean_formal}), and
$\|H_{\mathrm{local}} - D_W^L\|_F^2 = \kappa_g^2 S(n_{\max})/(4\pi^2) +
D(n_{\max})$ with $S, D$ polynomials in $n_{\max}$
(Corollary~\ref{cor:pythagorean_orthogonality}).  The sole transcendental,
$1/\pi^2$, enters through the posited $\beta = 2\pi/\kappa_g$ --- the
continuum KMS/Bisognano--Wichmann period imported as a normalisation (a
temporal-compactification $2\pi$ in the sense of Paper~35~\cite{paper35})
--- and is an input, not an output of the finite construction.  The
orthogonality is a statement about two chosen operators;\ it does not show
that $D_W$ is the ``wrong'' modular generator, which only O0 could decide.
The $S^3$-specificity of the construction --- the Bargmann--Segal $S^5$
Hardy sector lacks a spinor sector and a half-integer wedge --- is the
fourth layer of the Paper~24~\S~5.3~\cite{paper24} Coulomb/HO asymmetry.

\textbf{(O4) Higher-rank compact non-abelian extensions.}
Paper~40~\cite{paper40_unified} extends the Paper~38 state-space
Gromov--Hausdorff convergence theorem to compact connected semisimple Lie
groups with bi-invariant metric (a circle factor carries the circle constant). The natural extension of the present paper would
verify that the geometric BW-$\alpha$ generator analog on higher-rank
groups (e.g., $K_\alpha^W = J_z$ with $J_z$ the Cartan generator
along the wedge-defining direction) still has integer spectrum on the
spinor harmonic basis.  Since any integer-weight Cartan generator
closes automatically, that extension would be as elementary as the
present one;\ the substantive higher-rank question is the analog of O0.

\textbf{(O5) Cross-manifold modular structures.}
The Paper~24~\cite{paper24} cross-manifold blocker
$\Tcal_{\sthree} \otimes \Tcal_{\mathrm{Hardy}(S^5)}$ remains
unaddressed by the present paper. A cross-manifold modular construction
(modular Hamiltonian on a tensor product of a round $\sthree$
Camporesi--Higuchi triple and a holomorphic-sector $S^5$
Bargmann--Segal triple) would require an analog of the Latr\'emoli\`ere
tensor-product propinquity of Paper~39~\cite{paper39} extended to
mixed Riemannian / Hardy-sector spectral triples. This is also open.

\paragraph{Final remark on the M1 master Mellin engine.}
An earlier version of this paper closed by identifying the $2\pi$ of
the modular period with the $2\pi$ of the M1 volume content of the
Paper~38 state-space-GH rate constant, ``at the operator-system level
at every finite cutoff''.  That identification is withdrawn
(2026-09-28).  The period of the present construction is $2\pi$
because the posited generator has integer spectrum;\ that is the
generic fact $e^{2\pi i n} = 1$, which says nothing about the
Paper~38 rate constant or about Bisognano--Wichmann.  The M1 reading
of the four-witness $2\pi$ remains at the metric-functional level
(Paper~32~\S~VIII), where it was before.

% =====================================================================
\section{Krein-carrier extension at finite cutoff}
\label{sec:lorentzian_closure}
% =====================================================================

\paragraph{Provenance.}
This section is the successor of the former standalone Paper~43
(\emph{Lorentzian extension of the four-witness Wick-rotation theorem on
truncated $\sthree \times \R$ spectral triples at finite cutoff}),
archived on 2026-09-28.  Its surviving content --- the Krein carrier and
Lorentzian Dirac (\S\ref{sec:krein_carrier}), the axiom audit at
Bizi--Brouder--Besnard signature $(4, 6)$ (\S\ref{sec:bbb_audit}), the
Lorentzian wedge (\S\ref{sec:krein_wedge}), the Krein-carrier period
closure and the six-witness statement (\S\ref{sec:krein_four_witness_theorem}),
the Riemannian-limit recovery and the signature-independence of the
$H_{\mathrm{local}}$ residual
(\S\S\ref{sec:riemannian_limit_identification}--\ref{sec:h_local_signature_independent}),
the Hilbert--Schmidt orthogonality (\S\ref{sec:pythagorean_orthogonality})
and the M3 convention finding (\S\ref{sec:m3_convention}) --- is collected
here with its current tiers.  Two headline readings of that paper are
withdrawn and appear here only as history:\ (i)~that the Krein-carrier
period closure is an operator-system analog of the four-witness theorem
at signature $(3, 1)$ --- it is the automatic integer-spectrum closure of
Theorem~\ref{thm:bw_alpha}, and signature-blind;\ (ii)~that the $1/\pi^2$
of the Hilbert--Schmidt residual is a master-Mellin-engine M1 output ---
it is the square of the posited normalisation $\beta = 2\pi/\kappa_g$.

Open question (O1) of \S\ref{sec:conclusion} named the extension to
signature $(3, 1)$ as the next step.  Sprint L2 (May~2026) carries the
Riemannian construction of
\S\S\ref{sec:setup}--\ref{sec:four_witness_collapse} onto a Krein
carrier at finite cutoff, via the van~den~Dungen~2016~\cite{vandungen2016}
Proposition~4.1 lift and the
Bizi--Brouder--Besnard~2018~\cite{bizi_brouder_besnard2018}
classification.  The construction is structurally additive:\ at
$\Nt = 1$ every object reduces bit-identically to its Riemannian
counterpart (Proposition~\ref{prop:riemannian_limit_l2}).  Numerical
statements below refer to the tested panel
$\nmax \le 3$, $\Nt \in \{1, 11, 21\}$ (supporting modules
\texttt{geovac/krein\_space\_construction.py},
\texttt{geovac/lorentzian\_dirac.py},
\texttt{geovac/connes\_axiom\_audit\_31.py},
\texttt{geovac/modular\_hamiltonian\_lorentzian.py}, with the
corresponding \texttt{tests/test\_*.py} files).

\subsection{The Krein carrier and the Lorentzian Dirac}
\label{sec:krein_carrier}
\label{sec:lorentzian_closure_architecture}

\paragraph{Convention lock.}
We use the West-coast metric $\eta = \mathrm{diag}(+1, -1, -1, -1)$ on
$\sthree \times \R$ and the Peskin--Schroeder chiral basis
\begin{equation}\label{eq:krein_gamma_chiral}
   \gamma^0 = \begin{pmatrix} 0 & I_2 \\ I_2 & 0 \end{pmatrix},
   \qquad
   \gamma^i = \begin{pmatrix} 0 & \sigma^i \\ -\sigma^i & 0 \end{pmatrix},
   \qquad
   \gamma^5 := i\gamma^0\gamma^1\gamma^2\gamma^3
   = \mathrm{diag}(-I_2, +I_2).
\end{equation}
The spatial space is the chirality-doubled Camporesi--Higuchi spinor
space $\Hilb_{\mathrm{GV}}^{\nmax}$ of dimension
$\tfrac{2}{3}\nmax(\nmax+1)(\nmax+2) = 4, 16, 40$ at $\nmax = 1, 2, 3$,
with basis labels $(n_{\mathrm{fock}}, l, m_j, \chi)$, $j = l + \tfrac12$,
$\chi = \pm 1$, on which the truthful Camporesi--Higuchi Dirac acts
diagonally,
\begin{equation}\label{eq:krein_DGV}
   \DGV\, \psi_{n_{\mathrm{fock}}, l, m_j, \chi}
   \;=\; \chi\,(n_{\mathrm{fock}} + \tfrac12)\,
   \psi_{n_{\mathrm{fock}}, l, m_j, \chi}.
\end{equation}
The chirality label is identified with the $\gamma^5$ eigenvalue (up to a
global sign):\ $\gamma^5_{\mathrm{spatial}}\psi_{\cdot,\chi} = -\chi\,\psi_{\cdot,\chi}$.
The temporal slot is the bounded uniform grid
$t_k = -T_{\max} + k\cdot 2T_{\max}/(\Nt - 1)$, $k = 0, \ldots, \Nt - 1$
(the singleton $\{0\}$ at $\Nt = 1$).  It is \emph{not} compactified to
$S^1_\beta$, which would create closed timelike curves on
$\sthree \times S^1_\beta$ and conflict with the spacelike-reflection
ingredient of the van~den~Dungen lift.

\begin{definition}[Krein carrier]\label{def:krein_space}
$\Krein_{\nmax, \Nt} := \Hilb_{\mathrm{GV}}^{\nmax} \otimes \C^{\Nt}$,
with fundamental symmetry $\mathcal{J} := \gamma^0_{\mathrm{spatial}}
\otimes I_{\Nt}$ (the chirality swap) and Krein form
$\langle \psi, \phi \rangle_{\Krein} := \langle \psi, \mathcal{J}\phi\rangle$.
\end{definition}

\begin{proposition}[Krein axioms, {[SYMBOLIC]}]\label{prop:krein_axioms}
$\mathcal{J}^2 = I$, $\mathcal{J}^* = \mathcal{J}$, and
$\Krein = \Krein^+ \oplus \Krein^-$ orthogonally with
$\dim\Krein^\pm = \dim\Krein/2$.  In floating point all residuals are
exactly $0$ on the panel, because $\mathcal{J}$ is a real $\{0,1\}$
permutation matrix tensored with the identity.
\end{proposition}

\begin{definition}[Lorentzian Dirac]\label{def:lorentzian_dirac}
\begin{equation}\label{eq:DL_def}
   \DL := i\bigl[\gamma^0 \otimes \partial_t + \DGV \otimes I_{\Nt}\bigr],
   \qquad
   (\partial_t)_{k,l} = \frac{\delta_{l,k+1} - \delta_{l,k-1}}{2\Delta_t},
   \quad \Delta_t = \frac{2T_{\max}}{\Nt - 1},
\end{equation}
with Dirichlet zero boundary conditions, so that $\partial_t$ is real
skew-symmetric.  The prefactor $i^t = i$ is the $(s, t) = (3, 1)$
instance of van~den~Dungen's Proposition~4.1.
\end{definition}

\begin{proposition}[Krein-self-adjointness, {[SYMBOLIC]}]
\label{prop:krein_self_adjoint}
$\DL^\times := \mathcal{J}\DL^\dagger\mathcal{J} = \DL$.
\end{proposition}

\begin{proof}
$\DL^\dagger = -i(-\gamma^0\otimes\partial_t + \DGV\otimes I)$, using
$\partial_t^\dagger = -\partial_t$ and $\DGV^\dagger = \DGV$.  Conjugating
by $\mathcal{J}$ and using $(\gamma^0)^2 = I$ together with
$\gamma^0\DGV\gamma^0 = -\DGV$ (the chirality swap reverses the sign of
the chirality-diagonal eigenvalue~\eqref{eq:krein_DGV}) gives
$i(\gamma^0\otimes\partial_t + \DGV\otimes I) = \DL$.  The other
prefactor $-i$ would give $\DL^\times = -\DL$.
\end{proof}

At $\Nt = 1$ the grid is a point, $\partial_t = 0$, and
\begin{equation}\label{eq:DL_riemannian_limit}
   \DL\big|_{\Nt = 1} = i\,\DGV,
\end{equation}
the Riemannian Camporesi--Higuchi Dirac up to the global factor $i$
that mediates between Krein- and Hilbert-self-adjointness.  The chirality
grading on $\Krein$ is $\gamma^5_{\Krein} := \gamma^5_{\mathrm{spatial}}
\otimes I_{\Nt}$.

\subsection{Axiom audit at Bizi--Brouder--Besnard signature $(4, 6)$}
\label{sec:bbb_audit}

The Bizi--Brouder--Besnard classification of indefinite spectral triples
by $(m, n) = (t + s, t - s) \bmod 8$ (their p.~3) places the Lorentzian
signature with one time and three space dimensions, $(t, s) = (1, 3)$, at
$(m, n) = (4, -2 \bmod 8) = (4, 6)$ (their Table~3 entry $(4, 6)$), with
primitive signs
$\varepsilon = +1$, $\varepsilon'' = -1$, $\kappa = -1$, $\kappa'' = +1$
(their Table~1 and eq.~(3)).  The real structure is the lift of the $\Cl(3, 1)$
charge conjugation $U_4 = i\gamma^2 = (i\sigma_y)_{\mathrm{chir}}
\otimes (i\sigma^2)_{\mathrm{spin}}$,
\begin{equation}\label{eq:JL_def}
   \JL := U_L K, \qquad U_L := U_{L,\mathrm{spatial}} \otimes I_{\Nt},
   \qquad
   U_{L,\mathrm{spatial}}\,\psi_{n_{\mathrm{fock}}, l, m_j, \chi}
   = \sigma_{\mathrm{chir}}(\chi)\,(-1)^{(2l+1-2m_j)/2}\,
     \psi_{n_{\mathrm{fock}}, l, -m_j, -\chi},
\end{equation}
with $K$ complex conjugation and $\sigma_{\mathrm{chir}}(\pm 1) = \mp 1$.

\begin{theorem}[BBB-predicted-sign relations, {[SYMBOLIC]}]
\label{thm:connes_audit}
On $\Krein_{\nmax, \Nt}$:\ (i)~$\JL^2 = +I$;\
(ii)~$\{\JL, \gamma^5_{\Krein}\} = 0$;\
(iii)~$\{\JL, \mathcal{J}\} = 0$;\
(iv)~$\JL\DL = \DL\JL$.
All four residuals are exactly $0$ on the panel.
\end{theorem}

\begin{proof}
(i)~$(i\sigma_y)^2 = -I_{\mathrm{chir}}$ and $(i\sigma^2)^2 = -I_{\mathrm{spin}}$,
so $U_L^2 = +I$; $U_L$ is real, so $\JL^2 = U_L\overline{U_L} = +I$.
(ii),(iii)~$i\sigma_y$ anticommutes with both $\gamma^5$ (diagonal
$\mp$ in the chirality factor) and the chirality swap $\gamma^0$.
(iv)~$\gamma^0$, $\partial_t$ and $\DGV$ are real, so $K$ flips the sign
of the prefactor $i$;\ conjugation by $U_L$ flips the sign of $\gamma^0$
(anticommutation) and of $\DGV$ ($U_L$ reverses $\chi$
in~\eqref{eq:krein_DGV}).  The three signs cancel:\
$\JL\DL\JL^{-1} = \DL$.
\end{proof}

\begin{remark}[Order-zero and order-one conditions]
\label{rem:order_zero_one}
The order-zero condition $[a, \JL b\JL^{-1}] = 0$ and the order-one
condition $[[\DL, a], \JL b\JL^{-1}] = 0$, for multipliers $a, b$ of the
operator system, fail at finite cutoff, with Frobenius residuals up to
$\sim 7\%$ and $\sim 11\%$ on a sample of three multipliers per panel
cell.  These are \emph{not} finite-resolution artifacts expected to
vanish:\ on the Riemannian side, Paper~32~\cite{paper32} (\S~IV) measures
the order-zero failure \emph{growing} with the cutoff over
$\nmax = 2, 3, 4$ and, on a fixed generator set, approaching a nonzero
limit.  The cutoff dependence of the Lorentzian residuals has not been
measured.  (The former Paper~43 called them artifacts ``expected to vanish
in the GH limit''; withdrawn 2026-09-28.)
\end{remark}

\begin{proposition}[Universal anticommutation fails on truthful $\DGV$,
{[SYMBOLIC]}]\label{prop:bbb_axiom_failure}
The BBB universal relation $\{\chi, D\} = 0$ (their \S~5(v), required
at every signature) fails:\
$\{\gamma^5_{\Krein}, \DL\} = 2i\,(\gamma^5\DGV)\otimes I_{\Nt} \ne 0$,
with Frobenius norm $2\Frob{\DGV} = 6.00,\ 18.33,\ 38.88$ at
$\Nt = 1$, $\nmax = 1, 2, 3$.
\end{proposition}

\begin{proof}
$\{\gamma^5, \gamma^0\} = 0$ kills the temporal piece;\ $\gamma^5$ and
$\DGV$ are both chirality-diagonal, so they commute and
$\{\gamma^5, \DGV\} = 2\gamma^5\DGV$.
\end{proof}

\paragraph{A mutually inconsistent triple.}
\label{sec:bbb_axiom_finding}
(P1)~chirality identified with the $\gamma^5$ eigenvalue;\
(P2)~the Camporesi--Higuchi $\DGV$ chirality-diagonal~\eqref{eq:krein_DGV};\
(P3)~the BBB universal axiom $\{\chi, D\} = 0$.
Any two are compatible;\ (P1)$+$(P2) force $[\gamma^5, \DGV] = 0$,
contradicting (P3).  Three resolutions exist:\ (R1)~keep truthful $\DGV$
and record the failure as scope;\ (R2)~use the chirality-flipping
off-diagonal variant $\DGV^{\mathrm{offdiag}}$, which anticommutes with
$\gamma^5$ but breaks bit-exact Riemannian-limit recovery;\ (R3)~redefine
$\gamma^5$ as an off-diagonal grading, abandoning (P1).  This section
adopts (R1):\ the construction is a Krein-self-adjoint pair with the four
BBB-predicted-sign relations of Theorem~\ref{thm:connes_audit}, but not a
full BBB indefinite spectral triple in the strict sense of their
\S~5(v).  The finding is new at $(3, 1)$:\ the Riemannian triple has
odd KO-dimension and no chirality axiom to fail.  It does not affect the
period closures below, which depend only on the generator
$\KL^{\alpha, W}$, not on $\DL$.

\subsection{Lorentzian wedge, Gibbs state and generator}
\label{sec:krein_wedge}

\begin{definition}[Lorentzian hemispheric wedge]\label{def:wedge_L}
$P_{W_L} := P_W^{\mathrm{spatial}} \otimes P_{t \ge 0}$, with
$P_W^{\mathrm{spatial}}$ the hemispheric wedge of
Definition~\ref{def:wedge} and $P_{t\ge 0}$ the diagonal projector on
$\C^{\Nt}$ selecting $t_k \ge 0$.  Then $P_{W_L}$ is a Hermitian
idempotent, $\dim W_L = \tfrac12\dim\Hilb_{\mathrm{GV}}^{\nmax}\cdot
|\{k : t_k \ge 0\}|$, and $P_{W_L}|_{\Nt=1} = P_W^{\mathrm{spatial}}$.
\end{definition}

The generator is $\KL^\alpha := K_\alpha \otimes I_{\Nt}$ with
$K_\alpha$ of Definition~\ref{def:K_alpha_W}; its wedge restriction
$\KL^{\alpha, W}$ is diagonal with the positive odd integers
$\mathrm{two}\_m_j$ as eigenvalues.  Under the Bisognano--Wichmann choice
$H_{\mathrm{local}} := \KL^{\alpha, W}/\beta$ the wedge Gibbs state
\begin{equation}\label{eq:rho_WL}
   \rho_W^L := e^{-\beta H_{\mathrm{local}}}/Z = e^{-\KL^{\alpha, W}}/Z
\end{equation}
is $\beta$-independent.  Its GNS Hilbert--Schmidt space is
$M_{\dim W_L}(\C)$, with modular operator
$\Delta_L(X) = \rho_W^L X (\rho_W^L)^{-1}$, Tomita conjugation
$\JL^{\TT}(X) = (\rho_W^L)^{1/2} X^* (\rho_W^L)^{-1/2}$ satisfying
$(\JL^{\TT})^2 = +I$, and modular Hamiltonian
$\KL^{\TT} := -\log\Delta_L$, whose spectrum is
$\{n_j - n_k\} \subset \Z$ for $\{n_j\} = \mathrm{Spec}\,\KL^{\alpha, W}$.
Three distinct antilinear operators are now in play and must not be
conflated:\ $J_{\mathrm{GV}}$ on $\Hilb_{\mathrm{GV}}$
($J_{\mathrm{GV}}^2 = -I$), the Krein real structure $\JL$ on $\Krein$
($\JL^2 = +I$), and the state-dependent Tomita conjugation $\JL^{\TT}$
on the GNS space ($(\JL^{\TT})^2 = +I$).

\subsection{Krein-carrier period closure at finite cutoff}
\label{sec:krein_four_witness_theorem}

\begin{theorem}[Krein-carrier period closure, {[SYMBOLIC]}; automatic,
as Theorem~\ref{thm:bw_alpha}]
\label{thm:krein_four_witness}
With $\sigma_t^{L,\alpha}(O) := e^{it\KL^{\alpha, W}} O e^{-it\KL^{\alpha, W}}$
and $\sigma_t^{L,\TT}(a) := \Delta_L^{it}(a)$ on
$\Bcal(\Krein_{W_L})$, for every $(\nmax, \Nt)$:
\begin{enumerate}\setlength{\itemsep}{2pt}
\item[(a)] $\sigma_{2\pi}^{L,\alpha}(O) = O$ for every $O$;
\item[(b)] $\sigma_{2\pi}^{L,\TT}(a) = a$ for every $a$;
\item[(c)] $\sigma_t^{L,\TT}(a) = \sigma_{-t}^{L,\alpha}(a)$ for every
$t \in \R$;
\item[(d)] \emph{(six-witness statement, by construction.)}\ the six
instantiations $(\BW, \HH_{M=1}, \HH_{M=2}, \Sew_{M=1}, \Unruh_{a=1},
\Unruh_{a=2})$ yield identical $\rho_W^L$, $\Delta_L$ and $\KL^{\TT}$.
\end{enumerate}
In floating point, the residuals of (a)--(c) sit at the
$10^{-16}$--$10^{-15}$ round-off level on the panel, and the
cross-witness residual of (d) is exactly $0$.
\end{theorem}

\begin{proof}
$\KL^{\alpha, W}$ is diagonal with integer eigenvalues, so
$e^{\pm 2\pi i\KL^{\alpha, W}} = I$, which gives (a).  Since
$\rho_W^L = e^{-\KL^{\alpha, W}}/Z$, $\Delta_L^{it}(a) =
e^{-it\KL^{\alpha, W}} a\, e^{it\KL^{\alpha, W}}$ (the $Z$ cancels),
which is (c), and (b) follows from (a) and (c).  For (d), the witnesses
differ only in $\kappa_g$, which cancels in
$\beta H_{\mathrm{local}} = \KL^{\alpha, W}$.
\end{proof}

Every clause is elementary and none identifies a Bisognano--Wichmann
flow.  (a) and (b) hold for \emph{any} integer-spectrum generator and
\emph{any} operator;\ (c) is Tomita's theorem for a Gibbs state built
from the generator, so it is circular as a check on the generator;\ (d)
holds because the six witnesses share one generator.  The closure is
also \emph{signature-blind}:\ the posited generator is
compact (integer spectrum, [SYMBOLIC]), so its flow is a $\beta = 2\pi$
circle rather than a non-compact boost, and on the argued reading
[ARGUED] the $(3, 1)$ label is a carrier choice the closure cannot
distinguish from the Euclidean one (WH7;\ this compactness argument is
separate from Paper~45's $K^+$ annihilation theorem~\cite{paper45},
which concerns the Krein-positive compression).  The former Paper~43 read (a)--(d)
as an operator-system analog of the four-witness theorem at signature
$(3, 1)$;\ that reading is withdrawn (2026-09-28).  The finite-cutoff
test that would bear on it is the Riemannian O0 of \S\ref{sec:conclusion}.

\subsection{Riemannian-limit recovery at $\Nt = 1$}
\label{sec:riemannian_limit_identification}

\begin{proposition}[Riemannian-limit recovery, {[SYMBOLIC]}]
\label{prop:riemannian_limit_l2}
At $\Nt = 1$ every construction of this section reduces bit-identically
(Frobenius residual exactly $0$) to its Riemannian counterpart:\
$P_{W_L} \to P_W$ (Definition~\ref{def:wedge}),
$\KL^\alpha \to K_\alpha$ (Definition~\ref{def:K_alpha_W}),
$\KL^{\alpha, W} \to K_\alpha^W$, $\rho_W^L \to \rho_W$
(\S\ref{sec:kms_state}), $\Delta_L \to \Delta$ and
$\KL^{\TT} \to K_{\TT}$ (\S\ref{sec:tt_construction}), and
$\DL \to i\DGV$~\eqref{eq:DL_riemannian_limit}.
\end{proposition}

\begin{proof}
At $\Nt = 1$, $\partial_t = 0$, $P_{t\ge0} = I_1$, and every Kronecker
product with $I_1$ is the identity operation.
\end{proof}

This is the falsifier that makes the construction a lift of the
Riemannian one rather than a parallel construction; backed by
\texttt{tests/test\_modular\_hamiltonian\_lorentzian.py}.  At $\Nt > 1$
the temporal slot adds content (the temporal derivative in $\DL$, the
larger wedge) invisible at $\Nt = 1$.

\subsection{Signature-independence of the $H_{\mathrm{local}}$ residual}
\label{sec:h_local_signature_independent}

\begin{theorem}[$H_{\mathrm{local}}$ residual at the Riemannian limit,
{[SYMBOLIC]}]\label{thm:h_local_sig_independent}
At $\Nt = 1$,
$\Frob{H_{\mathrm{local}} - \DL^W} = \Frob{H_{\mathrm{local}} - D_W^{\mathrm{GV}}}$,
with values $2.1332$, $6.5275$, $13.854$ at $\nmax = 1, 2, 3$ for
$\kappa_g = 1$ ($\beta = 2\pi$).
\end{theorem}

\begin{proof}
$\DL^W|_{\Nt=1} = i D_W^{\mathrm{GV}}$ and $H_{\mathrm{local}}$ is real,
so the cross term $2\,\mathrm{Re}\,\Tr(H_{\mathrm{local}}\, iD_W^{\mathrm{GV}})$
of the Lorentzian side vanishes (the trace of a product of two real
matrices is real), and the Riemannian cross term
$2\,\Tr(H_{\mathrm{local}} D_W^{\mathrm{GV}})$ vanishes by
Theorem~\ref{thm:pythagorean_formal} below.  Both sides equal
$\Frob{H_{\mathrm{local}}}^2 + \Frob{D_W^{\mathrm{GV}}}^2$.
\end{proof}

The \S\ref{sec:derived_hamiltonian} / O3 distinction between the wedge
Dirac and the posited generator therefore persists verbatim at
signature $(3, 1)$ at the Riemannian limit, for the elementary reason in
the proof.  At $\Nt > 1$ the Lorentzian residual grows with the temporal
resolution [MEASURED;\ a PM measurement of 2026-09-28 at
$\nmax = 1, 2, 3$, $\Nt \le 41$, $T_{\max} = 1$, outside the declared test
panel and not test-backed]:\ at $\kappa_g = 1$ the finite-difference
$\partial_t$ term comes to dominate and the growth approaches $\Nt^{3/2}$
(residual ratio $\approx 2.8$ per doubling at $\nmax = 1$, lower at
$\nmax = 2, 3$ where the spatial part is larger) --- from the temporal-derivative content of $\DL$, which
$\KL^{\alpha, W}/\beta$ cannot capture.  This is a property of the posited
construction;\ it does not show which generator a Camporesi--Higuchi
wedge state selects (O0, open).

\subsection{Hilbert--Schmidt orthogonality:\ a chirality-parity identity}
\label{sec:pythagorean_orthogonality}

Let $\Pi_W$ be the wedge chirality grading ($\gamma^5_{\Krein}$ restricted
to the wedge;\ $\Pi_W^2 = I$, $\Pi_W^\dagger = \Pi_W$, $\Tr\Pi_W = 0$).

\begin{lemma}[Parity]\label{lem:parity_HS}
If $\Pi$ is a unitary involution, $[\Pi, A] = 0$ and $\{\Pi, B\} = 0$,
then $\Tr(A^\dagger B) = 0$.
\end{lemma}

\begin{proof}
$\Tr(A^\dagger B) = \Tr(\Pi A^\dagger\Pi\,\Pi B\Pi) = -\Tr(A^\dagger B)$.
\end{proof}

\begin{lemma}[Chirality pairing]\label{lem:pairing_HS}
Let $\Pi = \mathrm{diag}(\chi_i)$ with $\pm1$ eigenspaces paired
group-by-group, let $A$ be diagonal with equal entries on each pair, and
let $B = \Pi|B|$ with $|B|$ diagonal and equal on each pair.  Then
$\Tr(A^\dagger B) = 0$.
\end{lemma}

\begin{proof}
Each pair contributes $\bar a\,(b_+ + b_-) = \bar a\,(|b| - |b|) = 0$.
\end{proof}

\begin{theorem}[Hilbert--Schmidt orthogonality, {[SYMBOLIC]}]
\label{thm:pythagorean_formal}
For every $(\nmax, \Nt)$, every rescaling of the generator (in
particular every $\beta$ and every witness $\kappa_g$),
\[
   \big\langle H_{\mathrm{local}}, \DL^W\big\rangle_{\mathrm{HS}}
   := \Tr\big(H_{\mathrm{local}}^\dagger\DL^W\big) = 0.
\]
\end{theorem}

\begin{proof}
$H_{\mathrm{local}} = \KL^{\alpha, W}/\beta$ is diagonal with entries
$\mathrm{two}\_m_j/\beta$ independent of $\chi$, so it is $\Pi_W$-even and
$\chi$-independent.  Split $\DL^W$ into its temporal part
$i\gamma^0\otimes\partial_t$, which is $\Pi_W$-odd because
$\{\gamma^5, \gamma^0\} = 0$ (Lemma~\ref{lem:parity_HS}), and its spatial
part $i\DGV^W\otimes I$, which is $\Pi_W$-\emph{even} but sign-paired,
$\DGV^W = \Pi_W|\DGV^W|$ with $|\DGV^W| = n_{\mathrm{fock}} + \tfrac12$
independent of $\chi$~\eqref{eq:krein_DGV} (Lemma~\ref{lem:pairing_HS}).
Both traces vanish.
\end{proof}

The spatial factor is not $\Pi_W$-odd:\ it commutes with $\Pi_W$, which is
precisely the failure of the BBB universal axiom
(Proposition~\ref{prop:bbb_axiom_failure}).  The orthogonality is a
chirality-parity identity about two chosen operators;\ it is independent
of $\beta$ and does not decide which generator a wedge state selects.

\begin{corollary}[Pythagorean residual]\label{cor:pythagorean_orthogonality}
At $\Nt = 1$, with $\beta = 2\pi/\kappa_g$,
\begin{equation}\label{eq:pythagorean_residual}
   \Frob{H_{\mathrm{local}} - \DL^W}^2
   = \frac{\kappa_g^2\,S(\nmax)}{4\pi^2} + D(\nmax),
   \quad
   S = \frac{\nmax(\nmax+1)(\nmax+2)(2\nmax^2+4\nmax-1)}{15},
   \quad
   D = \frac{\nmax(\nmax+1)(\nmax+2)(2\nmax+1)(2\nmax+3)}{20}.
\end{equation}
\end{corollary}

The corollary follows from Theorem~\ref{thm:pythagorean_formal} by
expanding the norm:\ $S = \Frob{\KL^{\alpha, W}}^2 =
4\sum_{\mathrm{wedge}} m_j^2$ is the cumulative wedge angular-momentum
trace and $D = \Frob{\DL^W}^2 = \sum_{\mathrm{wedge}}
(n_{\mathrm{fock}} + \tfrac12)^2$ the squared-Dirac trace.
The two polynomial closed forms are [MEASURED]:\ in the original
measurement they were identified at $\nmax = 1, \ldots, 6$ (PSLQ at 100
digits on $\{r^2, 1, 1/\pi^2\}$), and the orthogonality was measured
across a larger panel and across the six witnesses at $\nmax \le 3$ (max
$8.9\times10^{-16}$);\ neither the PSLQ step nor the six-witness check is
in the regression suite.  The backing test
\texttt{tests/test\_paper43\_pythagorean\_hs\_orthogonality.py} recomputes
the closed forms from the live wedge operators at $\nmax = 1, \ldots, 4$
($\kappa_g = 1$) and checks the orthogonality on 7 cells (Riemannian
$\nmax = 2, 3, 4$;\ Lorentzian $(\nmax, \Nt) = (2,1), (3,1), (2,3),
(3,3)$, tolerance $10^{-12}$).  The decomposition persists at $\Nt > 1$, the
temporal-derivative contribution adding only to the rational part.

\paragraph{Transcendental tag.}
The sole transcendental is $1/\pi^2$, and it enters only through the
posited normalisation $\beta = 2\pi/\kappa_g$ in
$H_{\mathrm{local}} = K_\alpha^W/\beta$, squared by the Frobenius norm.
In Paper~18/Paper~34~\cite{paper34} terms it is the $2\pi$ of a thermal
(temporal-compactification) circle (Paper~35~\cite{paper35}) imported as
input --- the continuum Bisognano--Wichmann/KMS period --- not an output
of any spectral computation on the truncated triple.  The residual lies
in $\Q[\nmax] \oplus \pi^{-2}\Q[\nmax]$, the algebraic class of the M1
sector of the master Mellin engine (Paper~32~\S~VIII), but no Mellin
mechanism producing it is exhibited, and it is not the M1 content of the
Paper~38 rate constant.  The former Paper~43 read it as the squared
Hopf-base measure $1/\Vol(S^1)^2$ and a $k = 0$ instance of the
case-exhaustion theorem;\ withdrawn (2026-09-28).

\subsection{The M3 sub-mechanism is sectional in $n_{\mathrm{fock}}$-parity}
\label{sec:m3_convention}

The Sprint~L0 prediction (\S\ref{sec:l0_audit}) that the M3
sub-mechanism of the master Mellin engine trivialises at $(3, 1)$ on
chirality-symmetric spectrum is convention-dependent.  Under the
$n_{\mathrm{fock}}$-parity convention of the Paper~28~\cite{paper28}
vertex-parity sum,
$\sum_{n=1}^{N} g_n(-1)^n/(n+\tfrac12)^4$ with $g_n = 2n(n+1)$ counting
both chiralities, the value at $(3, 1)$ is bit-identical to the
Riemannian one at $\nmax \le 5$, so M3 does not trivialise.  Under a
chirality-pairing convention ($\chi = +1$ even, $\chi = -1$ odd) the
difference vanishes tautologically, but that convention does not access
the Catalan / Dirichlet-$\beta$ content of Paper~28.  The M3 sub-mechanism
is therefore sectional in $n_{\mathrm{fock}}$-parity, not in spacetime
signature;\ the case-exhaustion classification of Paper~32~\S~VIII is
unchanged.

\subsection{Honest scope at finite cutoff}
\label{sec:lorentzian_closure_scope}

\begin{itemize}\setlength{\itemsep}{2pt}
\item Everything in this section is a finite-$(\nmax, \Nt)$ statement.
No continuum limit is taken.  A Lorentzian analog of the
Latr\'emoli\`ere propinquity~\cite{latremoliere2018,
latremoliere_metric_st_2017} is absent from the literature
(Toyota~\cite{toyota2023}, Farsi--Latr\'emoli\`ere~\cite{farsi_latremoliere2024,
farsi_latremoliere2025}, Hekkelman--McDonald~\cite{hekkelman_mcdonald2024b}
are Riemannian), and the framework's attempt was descoped on 2026-06-09
(Paper~45's $K^+$ annihilation theorem:\ the natural Krein-positive
compression fails structurally).  Separately, the posited
generator is compact [SYMBOLIC];\ that no Lorentzian metric statement is
therefore available at finite cutoff for it is [ARGUED] (Paper~45 Q1).  The
Nieuviarts twist morphism~\cite{nieuviarts2025a, nieuviarts2025b_proceedings}
is restricted to even dimension and does not apply to $\sthree$.
\item The construction uses truthful $\DGV$ (resolution R1 of
\S\ref{sec:bbb_axiom_finding}), so it is not a full BBB indefinite
spectral triple;\ the period closures are $\DL$-independent and
unaffected.
\item The cross-manifold product with the Bargmann--Segal $S^5$ Hardy
sector (Paper~24~\cite{paper24}~\S~V) and the inner-factor calibration
question are not addressed.
\item The operator-system substrate on the full Krein space (the
propagation number and its envelope dependence) is
Paper~44~\cite{paper44}, which takes the carrier of
\S\ref{sec:krein_carrier} as input.
\end{itemize}

% =====================================================================
\appendix
\section{Computational verification protocol}
\label{app:computational_verification}
% =====================================================================

The numerical verification of the period closures of
Theorems~\ref{thm:bw_alpha} and~\ref{thm:bw_gamma} and of the
cross-witness collapse of Table~\ref{tab:cross_witness} is
implemented in the supporting Python module
\verb|geovac/modular_hamiltonian.py| of the present author's framework
codebase ($\sim$1550 lines), with test coverage in
\verb|tests/test_modular_hamiltonian.py| ($\sim$960 lines, 67 tests:\
63 fast plus 4 slow). The structure is briefly described here.

\subsection{Module API}\label{sec:module_api}

The principal classes are:
\begin{itemize}\setlength{\itemsep}{2pt}
\item \texttt{HemisphericWedge(axis="hopf")}:\ implements the projection
$P_W$ of Definition~\ref{def:wedge} via the equatorial-reflection
involution $R_{\mathrm{polar}}$ on the full-Dirac spinor basis.
\item \texttt{TomitaConjugation(U)}:\ implements the antilinear
$J_{\TT}$ via a stored unitary $U$ with action
$J_{\TT}(\psi) = U \cdot \overline{\psi}$;\ exposes
\texttt{verify\_J\_squared\_positive\_identity} and
\texttt{verify\_J\_squared\_negative\_identity\_DOES\_NOT\_HOLD} as
explicit sanity checks distinguishing $J_{\TT}$ from
$J_{\mathrm{GV}}$.
\item \texttt{KMSState(beta, H)}:\ implements the wedge-restricted
Gibbs state $\omega_{\beta, H}$ of~\eqref{eq:gibbs_state_def},
including the canonical analytic-continuation verification of the
KMS condition at the expectation level.
\item \texttt{ModularHamiltonian(wedge, basis, D, kappa\_g)}:\
implements the BW-$\alpha$ construction with $K = K_\alpha^W$, the
real-time flow $\sigma_t^{\alpha}$, the period-closure verifier
\texttt{verify\_modular\_periodicity}, and accessors to the
Tomita--Takesaki structure (\texttt{tomita\_structure},
\texttt{k\_alpha\_geometric}, \texttt{k\_tomita},
\texttt{compare\_alpha\_vs\_tomita}).
\item \texttt{TomitaModularStructure(rho)}:\ implements the
BW-$\gamma$ Tomita construction on the GNS Hilbert--Schmidt space,
with $\Delta = (\rho^{-1})^T \otimes \rho$, $K_{\TT} = -\log \Delta$,
and the modular flow $\sigma_t^{\TT}(a) = \rho^{it} a \rho^{-it}$ on
the algebra.
\end{itemize}
Witness factories \texttt{for\_bisognano\_wichmann},
\texttt{for\_hartle\_hawking}, \texttt{for\_sewell},
\texttt{for\_unruh} instantiate
\texttt{ModularHamiltonian} with witness-specific $\kappa_g$ values:\
BW witness ($\kappa_g = 1$), HH ($\kappa_g = 1/(4M)$),
Unruh ($\kappa_g = a$). Each factory returns a
\texttt{ModularHamiltonian} object whose period-closure verifier
returns the label \texttt{STRONG\_IDENTIFICATION} at every tested
$\nmax$;\ the label records the automatic closure, not an
identification (Remark~\ref{rem:bw_alpha_automatic}).  The witness
factories share the generator $K_\alpha^W$ and differ only in
$\kappa_g$, so their agreement is by construction.

\subsection{Test coverage}\label{sec:test_coverage}

The test suite \verb|tests/test_modular_hamiltonian.py| verifies the
following:\ wedge projector idempotency and Hermiticity ($P_W^2 =
P_W$, $P_W^* = P_W$); Tomita conjugation properties ($J_{\TT}^2 = +I$,
$J_{\TT}^2 \ne -I$, antilinearity, operator-action consistency); KMS
state normalisation, partition-function positivity, real expectation
on Hermitian operators; modular periodicity $\sigma_{2\pi}(O) = O$ at
$\nmax \in \{2, 3, 4\}$ for the first five non-identity multipliers in
\verb|FullDiracTruncatedOperatorSystem|; cross-witness collapse at
the operator-action level; KMS expectation-level condition for sample
operator pairs $(A, B)$; propinquity rate cross-check
(residual $\ll \gamma_{\nmax}$); operator-system non-leakage
($\sigma_{2\pi}(M)$ stays in $\Op_{\nmax}$); modular flow group law
$\sigma_t \cdot \sigma_s = \sigma_{t+s}$ and identity at $t = 0$. The
slow tests verify the closure at $\nmax = 5$ ($\sim$80 seconds wall
time per slow test).

\subsection{Wall time and reproducibility}\label{sec:wall_time}

Per-$\nmax$ wall time on a standard laptop (Intel i7, Python 3.11,
NumPy 1.26):\ $\sim 0.4$ s at $\nmax = 2$, $\sim 3.5$ s at $\nmax = 3$,
$\sim 45$ s at $\nmax = 4$, $\sim 260$ s at $\nmax = 5$ for the
BW-$\alpha$ construction;\ $\sim 0.2$ s, $\sim 3.7$ s, $\sim 45$ s,
$\sim 545$ s for the BW-$\gamma$ construction (the $\nmax = 5$ cost
dominated by the $4900 \times 4900$ GNS-Hilbert--Schmidt
diagonalisation). Beyond $\nmax = 5$ the GNS construction becomes
expensive; the structural reading of the closure is unchanged at
higher $\nmax$ since the integer-spectrum mechanism holds at any
cutoff. Computational drivers and structured output JSON are
available in the present author's framework documentation under
\verb|debug/l1_modular_hamiltonian_compute.py| and
\verb|debug/l1_tighten_tomita_compute.py| respectively.

% =====================================================================
\section{Cross-references to the GeoVac framework}
\label{app:framework_crossref}
% =====================================================================

The present paper is part of an integrated NCG-research programme
developed by the present author. We briefly note the principal
cross-references for the reader's orientation.

\paragraph{Paper~38 (\cite{paper38}, state-space GH convergence on
$\sthree$).}
The convergence of the truncated Camporesi--Higuchi spectral triples
$\Tcal_{\nmax}$ to the continuum $\Tcal_{\sthree}$ in van~Suijlekom's
state-space Gromov--Hausdorff distance is established in Paper~38 with
explicit rate~\eqref{eq:paper38_main}. The L1' / L2 / L3 / L4
lemmas of Paper~38 supply the metric-side analytic infrastructure
that the present paper builds on;\ Paper~38's fifth lemma (the
propinquity assembly) was withdrawn on 2026-09-03 and its convergence
theorem is carried by the lifted-state construction, which the present
paper does not depend on either way. The state-space GH rate constant
$4/\pi = \Vol(S^2)/\pi^2$ carries the $k = 0$ M1 volume content;\ the
period closure of the present
Theorems~\ref{thm:bw_alpha}--\ref{thm:bw_gamma} is driven instead by
the integer spectrum of the posited generator, and no link between the
two is claimed (withdrawn 2026-09-28).

\paragraph{Paper~40 (\cite{paper40_unified}, universal $4/\pi$ on
compact semisimple Lie groups).}
The Paper~38 state-space Gromov--Hausdorff convergence extends to compact
connected semisimple Lie groups with bi-invariant metric (a circle factor
carries the circle constant), with the rate constant $4/\pi$ (dual-Coxeter normalisation of the $\SU(2)$ metric; $2/\pi$ on the unit round $S^3$, the constant scaling with the metric --- Paper~38, 2026-09-03)
reported universal across that class via the Plancherel-weight $\times$
Vandermonde-Jacobian cancellation (Paper~40~\S~3.2; rigorous at rank~1). The
modular structure of the present paper plausibly extends to the
higher-rank cases (Section~\ref{sec:conclusion} O4); the integer
spectrum of the posited rotation generator $K_\alpha^W$ on the spinor
harmonic basis is a property of $\SU(2)$ (where the half-integer
weights double to integers), and the higher-rank analog requires
verifying that the relevant Cartan generator along the wedge axis
still has the right spectral property.

\paragraph{Paper~32 \S VIII (master Mellin engine case-exhaustion
theorem).}
The case-exhaustion theorem~\cite[Thm.~thm:pi\_source\_case\_exhaustion]{paper32}
of Paper~32 \S VIII (in the present author's framework documentation)
states that every $\pi$ in any finite chain of Paper~34 projections
engages one of three master Mellin engine sub-mechanisms:\ M1 (Hopf-base
measure $\Vol(S^2)/\pi^2$ at operator order $k = 0$), M2 (Seeley--DeWitt
$\sqrt{\pi}\, \Q$ at $k = 2$), or M3 (vertex-parity Hurwitz / Catalan
$G$ / Dirichlet $L$ at $k = 1$). The $2\pi$ in the modular period of
the present paper comes from the integer spectrum of the posited
generator ($e^{2\pi i n} = 1$) and is not shown here to be the M1
signature.  The Bisognano--Wichmann reading of M1 on temporal
compactifications (\cite[Rem.~rem:bisognano\_wichmann\_reading]{paper32})
remains a structural correspondence at the metric-functional level;\
the present paper does not lift it (an earlier reading that it did is
withdrawn, 2026-09-28).

\paragraph{Paper~34 \S III.27 (Wick rotation projection).}
The Wick rotation projection in the present author's projection
taxonomy~\cite{paper34} is the structural mechanism through which the
framework's Euclidean spectral content is read as a Lorentzian
physical observable. The Sprint L0 audit
(\texttt{debug/lorentzian\_l0\_audit\_memo.md}) classifies this
projection in the WICK\_MAP\_FREE bucket at the metric-functional
level. The present paper does not change that status:\ its
finite-cutoff period closure is automatic for the posited generator,
so the framework still relies on the published Wick-rotation chain for
the four-witness period (withdrawn reading, 2026-09-28).

\paragraph{Paper~24 (\cite{paper24}, Coulomb / HO asymmetry).}
The cross-manifold blocker named in Paper~24 \S V (Riemannian $\sthree$
$\otimes$ holomorphic-sector $S^5$ Bargmann--Segal triple) carries
through to the modular construction:\ the present paper does not
address the cross-manifold modular structure, and the natural
extension would require a mixed Riemannian / Hardy-sector
tensor-product propinquity that does not yet exist in the literature.

\paragraph{Paper~2 (\cite{paper2}, fine-structure-constant
observation).}
The present paper does not impinge on Paper~2's status as an
Observation. The $\nmax = 3$ cutoff of
Table~\ref{tab:bw_alpha_residuals} happens to coincide with Paper~2's
selection-principle truncation ($\dim \Hilb_3 = 4 + 12 + 24 = 40$, the
cumulative node count, numerically equal to Paper~2's single-level
degeneracy $g_3^{\mathrm{Dirac}} = \Delta^{-1} = 40$ --- different objects
that agree only at $\nmax = 3$); this is a numerical coincidence at the
Hilbert-space dimension level, not a derivation of Paper~2's combination rule
$K = \pi(B + F - \Delta)$. The combination rule remains a numerical
observation without first-principles derivation.

% =====================================================================
\section*{Acknowledgments}
% =====================================================================

The author thanks Alain Connes, Carlo Rovelli, Walter van Suijlekom,
Fran\c{c}ois Latr\'emoli\`ere, Eva-Maria Hekkelman, Edward McDonald,
Malte Leimbach, Matilde Marcolli, Horacio Casini, Marina Huerta,
Robert Myers, and Rainer Verch for the foundational frameworks
on which this paper builds; remaining errors are the author's own.
This work was developed in an AI-augmented agentic workflow with
Anthropic's Claude:\ implementation of the supporting computational
modules, drafting of the internal proof memos used as supporting
documentation, and bibliographic collation were performed
collaboratively. All theorem statements, design choices, and
editorial decisions are the author's responsibility. The author
thanks the open-source NCG community (notably the SageMath, sympy,
mpmath, and scipy teams) for the computational infrastructure used in
the supporting numerical verifications.

% =====================================================================
\bibliographystyle{plainnat}
\begin{thebibliography}{99}

\bibitem[Bisognano \& Wichmann(1975)]{bisognano_wichmann1975}
J.~J.~Bisognano and E.~H.~Wichmann,
\newblock \emph{On the duality condition for a Hermitian scalar field},
\newblock J.~Math.~Phys. \textbf{16} (1975), 985--1007.

\bibitem[Bisognano \& Wichmann(1976)]{bisognano_wichmann1976}
J.~J.~Bisognano and E.~H.~Wichmann,
\newblock \emph{On the duality condition for quantum fields},
\newblock J.~Math.~Phys. \textbf{17} (1976), 303--321.

\bibitem[Bizi, Brouder \& Besnard(2018)]{bizi_brouder_besnard2018}
N.~Bizi, C.~Brouder, and F.~Besnard,
\newblock \emph{Space and time dimensions of algebras with
application to Lorentzian noncommutative geometry and quantum
electrodynamics},
\newblock J.~Math.~Phys. \textbf{59} (2018), 062303;
\newblock arXiv:1611.07062.

\bibitem[Camporesi \& Higuchi(1996)]{camporesi_higuchi1996}
R.~Camporesi and A.~Higuchi,
\newblock \emph{On the eigenfunctions of the Dirac operator on
spheres and real hyperbolic spaces},
\newblock J.~Geom.~Phys. \textbf{20} (1996), 1--18.

\bibitem[Casini \& Huerta(2009)]{casini_huerta2009}
H.~Casini and M.~Huerta,
\newblock \emph{Entanglement entropy in free quantum field theory},
\newblock J.~Phys. A: Math.~Theor. \textbf{42} (2009), 504007;
\newblock arXiv:0905.2562.

\bibitem[Casini, Huerta \& Myers(2011)]{casini_huerta_myers2011}
H.~Casini, M.~Huerta, and R.~C.~Myers,
\newblock \emph{Towards a derivation of holographic entanglement
entropy},
\newblock JHEP \textbf{05} (2011), 036;
\newblock arXiv:1102.0440.

\bibitem[Chamseddine \& Connes(2010)]{chamseddine_connes2010}
A.~H.~Chamseddine and A.~Connes,
\newblock \emph{Noncommutative geometry as a framework for
unification of all fundamental interactions including gravity. I},
\newblock Fortsch.~Phys. \textbf{58} (2010), 553--600.

\bibitem[Connes(1995)]{connes1995}
A.~Connes,
\newblock \emph{Noncommutative Geometry},
\newblock Academic Press, 1995.

\bibitem[Connes \& Rovelli(1994)]{connes_rovelli1994}
A.~Connes and C.~Rovelli,
\newblock \emph{Von Neumann algebra automorphisms and
time-thermodynamics relation in general covariant quantum theories},
\newblock Class.~Quantum Grav. \textbf{11} (1994), 2899--2917;
\newblock arXiv:gr-qc/9406019.

\bibitem[Connes \& van~Suijlekom(2021)]{connes_vs2021}
A.~Connes and W.~D.~van~Suijlekom,
\newblock \emph{Spectral truncations in noncommutative geometry and
operator systems},
\newblock Comm.~Math.~Phys. \textbf{383} (2021), 2021--2067;
\newblock arXiv:2004.14115.

\bibitem[Farsi \& Latr\'emoli\`ere(2024)]{farsi_latremoliere2024}
C.~Farsi and F.~Latr\'emoli\`ere,
\newblock \emph{Collapse in noncommutative geometry and spectral
continuity},
\newblock arXiv:2404.00240, 2024.

\bibitem[Farsi \& Latr\'emoli\`ere(2025)]{farsi_latremoliere2025}
C.~Farsi and F.~Latr\'emoli\`ere,
\newblock \emph{Continuity for the spectral propinquity of the
Dirac operators associated with an analytic path of Riemannian
metrics},
\newblock arXiv:2504.11715, 2025.

\bibitem[Hartle \& Hawking(1976)]{hartle_hawking1976}
J.~B.~Hartle and S.~W.~Hawking,
\newblock \emph{Path-integral derivation of black-hole radiance},
\newblock Phys.~Rev. D \textbf{13} (1976), 2188--2203.

\bibitem[Hekkelman(2022)]{hekkelman2022}
E.~Hekkelman,
\newblock \emph{Truncated geometry on the circle},
\newblock Lett.~Math.~Phys.\ \textbf{112} (2022), 20;
\newblock arXiv:2111.13865.

\bibitem[Hekkelman \& McDonald(2024b)]{hekkelman_mcdonald2024b}
E.~Hekkelman and E.~McDonald,
\newblock \emph{A noncommutative integral on spectrally truncated
spectral triples, and a link with quantum ergodicity},
\newblock J.~Funct.~Anal., to appear (2025);
\newblock arXiv:2412.00628.

\bibitem[Rieffel(2004)]{ucp_maps_2024}
M.~A.~Rieffel,
\newblock \emph{Matrix algebras converge to the sphere for quantum
Gromov--Hausdorff distance},
\newblock Mem. Amer. Math. Soc. \textbf{168} (2004), no.~796, 67--91;
arXiv:math/0108005.

\bibitem[Latr\'emoli\`ere(2016)]{latremoliere2018}
F.~Latr\'emoli\`ere,
\newblock \emph{The quantum Gromov--Hausdorff propinquity},
\newblock Trans.~Amer.~Math.~Soc. \textbf{368} (2016), 365--411.

\bibitem[Latr\'emoli\`ere(2017)]{latremoliere_metric_st_2017}
F.~Latr\'emoli\`ere,
\newblock \emph{The Gromov--Hausdorff propinquity for metric spectral
triples},
\newblock Adv.~Math. \textbf{404} (2022), Paper No.~108393, 56pp;
\newblock preprint arXiv:1811.10843, 2017.

\bibitem[Leimbach \& van~Suijlekom(2024)]{leimbach_vs2024}
M.~Leimbach and W.~D.~van~Suijlekom,
\newblock \emph{Gromov--Hausdorff convergence of spectral truncations
for tori},
\newblock Adv.~Math. \textbf{439} (2024), Paper No.~109496.

\bibitem[Marcolli \& van~Suijlekom(2014)]{marcolli_vs2014}
M.~Marcolli and W.~D.~van~Suijlekom,
\newblock \emph{Gauge networks in noncommutative geometry},
\newblock J.~Geom.~Phys. \textbf{75} (2014), 71--91;
\newblock arXiv:1301.3480.

\bibitem[Nieuviarts(2024)]{nieuviarts2024}
G.~Nieuviarts,
\newblock \emph{Signature change by a morphism of spectral
triples},
\newblock arXiv:2402.05839, 2024 (revised through v6, March 2026).

\bibitem[Nieuviarts(2025a)]{nieuviarts2025a}
G.~Nieuviarts,
\newblock \emph{Emergence of Lorentz symmetry from an
almost-commutative twisted spectral triple},
\newblock arXiv:2502.18105 (v3, May 2025).

\bibitem[Nieuviarts(2025b)]{nieuviarts2025b_proceedings}
G.~Nieuviarts,
\newblock \emph{Emergence of time from a twisted spectral triple
in almost-commutative geometry},
\newblock arXiv:2512.15450, December 2025 (revised May 2026).

\bibitem[Sewell(1982)]{sewell1982}
G.~L.~Sewell,
\newblock \emph{Quantum fields on manifolds: PCT and gravitationally
induced thermal states},
\newblock Ann.~Phys. (NY) \textbf{141} (1982), 201--224.

\bibitem[Strohmaier(2000)]{strohmaier2006}
A.~Strohmaier,
\newblock \emph{The Reeh--Schlieder property for quantum fields on
stationary spacetimes},
\newblock Comm.~Math.~Phys. \textbf{215} (2000), 105--118.

\bibitem[Takesaki(1970)]{takesaki1970}
M.~Takesaki,
\newblock \emph{Tomita's theory of modular Hilbert algebras and its
applications},
\newblock Lecture Notes in Mathematics 128, Springer, 1970.

\bibitem[Tomita(1967)]{tomita1967}
M.~Tomita,
\newblock \emph{Standard form of von Neumann algebras},
\newblock Lecture notes, RIMS Kokyuroku \textbf{14}, Kyoto University,
1967.

\bibitem[Toyota(2023)]{toyota2023}
R.~Toyota,
\newblock \emph{Quantum Gromov--Hausdorff convergence of spectral
truncations for groups with polynomial growth},
\newblock arXiv:2309.13469, 2023.

\bibitem[Unruh(1976)]{unruh1976}
W.~G.~Unruh,
\newblock \emph{Notes on black-hole evaporation},
\newblock Phys.~Rev. D \textbf{14} (1976), 870--892.

\bibitem[van~den~Dungen(2016)]{vandungen2016}
K.~van~den~Dungen,
\newblock \emph{Krein spectral triples and the fermionic action},
\newblock Math.~Phys.~Anal.~Geom. \textbf{19} (2016), 4;
\newblock arXiv:1505.01939.

\bibitem[Verch(2001)]{verch2001}
R.~Verch,
\newblock \emph{A spin-statistics theorem for quantum fields on
curved spacetime manifolds in a generally covariant framework},
\newblock Comm.~Math.~Phys. \textbf{223} (2001), 261--288.

\bibitem[Zhang \emph{et al.}(2020)]{zhu_casini2020}
J.~Zhang, P.~Calabrese, M.~Dalmonte, and M.~A.~Rajabpour,
\newblock \emph{Lattice Bisognano--Wichmann modular Hamiltonian in
critical quantum spin chains},
\newblock SciPost Phys.~Core \textbf{2} (2020), 007;
\newblock arXiv:2003.00315.

\bibitem[Loutey, internal Paper~2]{paper2}
J.~Loutey,
\newblock \emph{The Fine Structure Constant as a Spectral Coincidence
on the Hopf Fibration},
\newblock GeoVac internal preprint, papers/observations (2026),
DOI via Zenodo.

\bibitem[Loutey, internal Paper~24]{paper24}
J.~Loutey,
\newblock \emph{The Bargmann--Segal Lattice: $\pi$-Free Discretization
of the Harmonic Oscillator on $S^5$},
\newblock GeoVac internal preprint (2026), DOI via Zenodo.

\bibitem[Loutey, internal Paper~25]{paper25_hopf_gauge_structure}
J.~Loutey,
\newblock \emph{Hopf graph as a discrete lattice-gauge structure on
the Fock-projected $S^3$},
\newblock GeoVac internal preprint, papers/synthesis (2026),
DOI via Zenodo.

\bibitem[Loutey, internal Paper~28]{paper28}
J.~Loutey,
\newblock \emph{QED on $S^3$: Spectral Geometry of Perturbative
Transcendentals},
\newblock GeoVac internal preprint, papers/group5\_qed\_gauge (2026),
DOI via Zenodo.

\bibitem[Loutey, internal Paper~32]{paper32}
J.~Loutey,
\newblock \emph{The GeoVac Spectral Triple: A Synthesis Paper Locking
the Framework into the Marcolli--van~Suijlekom Gauge-Network Lineage},
\newblock GeoVac internal preprint, papers/synthesis (2026),
DOI via Zenodo.

\bibitem[Loutey, internal Paper~34]{paper34}
J.~Loutey,
\newblock \emph{GeoVac as a Two-Layer Framework: Projection Taxonomy
and the Empirical Matches Catalog},
\newblock GeoVac internal preprint, papers/observations (2026),
DOI via Zenodo.

\bibitem[Loutey, internal Paper~35]{paper35}
J.~Loutey,
\newblock \emph{Time as Projection:\ The Klein-Gordon Spectrum on
$S^3 \times \mathbb{R}$ and the Algebraic / Observation Split},
\newblock GeoVac internal preprint, papers/group6\_precision\_observations
(2026), DOI via Zenodo.

\bibitem[Loutey, internal Paper~38]{paper38}
J.~Loutey,
\newblock \emph{State-space Gromov--Hausdorff convergence of
truncated Camporesi--Higuchi spectral triples on $S^3$},
\newblock GeoVac internal preprint, papers/standalone (2026),
DOI via Zenodo.

\bibitem[Loutey, internal Paper~39]{paper39}
J.~Loutey,
\newblock \emph{Tensor-product propinquity convergence for
distinct-focal-length truncated Camporesi--Higuchi spectral triples
on $\sthree$},
\newblock GeoVac internal preprint, papers/standalone (2026),
DOI via Zenodo.

\bibitem[Loutey, internal Paper~40]{paper40_unified}
J.~Loutey,
\newblock \emph{State-space Gromov--Hausdorff convergence of spectral
truncations of compact semisimple Lie groups with bi-invariant metric},
\newblock GeoVac internal preprint, papers/standalone (2026),
DOI via Zenodo.

\bibitem[Loutey, internal Paper~44]{paper44}
J.~Loutey,
\newblock \emph{Operator-system extension of the BBB Krein spectral
triple at finite cutoff:\ propagation number and envelope dependence on
the Camporesi--Higuchi spinor bundle},
\newblock GeoVac internal preprint, papers/group1\_operator\_algebras
(2026), DOI via Zenodo.

\bibitem[Loutey, internal Paper~45]{paper45}
J.~Loutey,
\newblock \emph{Truncated Krein spectral triples over compact-temporal
$S^3 \times S^1_T$:\ the $K^{+}$ annihilation theorem},
\newblock GeoVac internal preprint, papers/group1\_operator\_algebras
(2026), DOI via Zenodo.

\end{thebibliography}

\end{document}
